\chapter{CLEBSCH-GORDAN COEFFICIENTS AND $3 j m$ SYMBOLS}\label{chap:8}\index{Clebsch-Gordan coefficient}\index{3jm symbol@$3jm$ symbol}
\section{DEFINITION}\label{sec:8.1}\index{vector-addition coefficients!definition}
The Clebsch-Gordan coefficients are vector addition coefficients.\index{vector-addition coefficients}\index{rotation group!irreducible representations of}\index{3nj symbols@$3nj$ symbols} They play an important role in the decomposition of reducible representations of rotation group into irreducible representations. All recoupling coefficients or $3 n j$ symbols can be determined as the sums of products of the Clebsch-Gordan coefficients. The aforesaid explains the extensive application of these coefficients in the quantum theory of angular momentum.

\subsection{The Clebsch-Gordan Coefficients}\label{sec:8.1.1}\index{Clebsch-Gordan coefficient!definition}
Let $\vect{j}_{1}$ and $\vect{j}_{2}$ be two angular momenta with projections $m_{1}$ and $m_{2}$ on the quantization axis. A ClebschGordan coefficients represents the probability amplitude that\index{probability amplitude}\isym{C-clebsch@$\clebsch{a}{\alpha}{b}{\beta}{c}{\gamma}$, Clebsch-Gordan coefficient} $\mathrm{j}_{1}$ and $\mathrm{j}_{2}$ are coupled into a resultant angular momentum $\vect{j}$ with projection $m$. In accordance with the vector addition rules $\vect{j}_{1}+\vect{j}_{2}=\vect{j}$, the Clebsch-Gordan coefficient vanishes unless the triangular conditions (triangular inequalities) are fulfilled\index{triangle condition}\index{triangular inequalities|see{triangle condition}}, i.e.,
\begin{equation}
\left|j_{1}-j_{2}\right| \leq j \leq j_{1}+j_{2}, \label{eq:8:1:1}
\end{equation}
and the requirement
\begin{equation}
m_{1}+m_{2}=m, \label{eq:8:1:2}
\end{equation}
is satisfied. It will also be assumed that arguments of the Clebsch-Gordan coefficients satisfy the following conditions:
\begin{enumerate}[label=(\alph*)]
\item $j_{1}, j_{2}, j$ are integer or half-integer non-negative numbers; \footnote{The extension of the Clebsch-Gordan coefficients to negative values of momenta will be considered in Sec.~\ref{sec:8.4.5}.
}
\item $m_{1}, m_{2}, m$ are integer or half-integer (positive or negative) numbers;
\item $\left|m_{1}\right| \leq j_{1},\left|m_{2}\right| \leq j_{2},|m| \leq j$;
\item $j_{1}+m_{1}, j_{2}+m_{2}, j+m, j_{1}+j_{2}+j$ are integer non-negative numbers.
\end{enumerate}
%no eq. 3 in book! conditions (a)-(d) above occupy equation number 8.1.3
\refstepcounter{equation}\label{eq:8:1:3}
The absolute value of a Clebsch-Gordan coefficient is given by
\begin{equation}
\left[\clebsch{j_{1}}{m_{1}}{j_{2}}{m_{2}}{j}{m}\right]^{2}=\frac{2 j+1}{8 \pi} \int_{0}^{2 \pi} d \alpha \int_{0}^{\pi} d \beta \sin \beta \int_{0}^{2 \pi} d \gamma D_{m_{1} m_{1}}^{j_{1}}(\alpha, \beta, \gamma) D_{m_{2} m_{2}}^{j_{2}}(\alpha, \beta, \gamma) D_{m m}^{j *}(\alpha, \beta, \gamma). \label{eq:8:1:4}
\end{equation}
The phase of the Clebsch-Gordan coefficients may be chosen in different ways. The phase convention proposed by Condon and Shortley \cite{ref010} is universally accepted.\index{Condon-Shortley phase convention}\index{phase!of the Clebsch-Gordan coefficients} In accordance with this convention the Clebsch-Gordan


coefficients are real and their phases are fixed by the additional relation
\begin{equation}
\clebsch{j_{1}}{m_{1}}{j_{2}}{m_{2}}{j}{m} \clebsch{j_{1}}{j_{1}}{j_{2}}{-j_{2}}{j}{j_{1}-j_{2}}=\frac{2 j+1}{8 \pi} \int_{0}^{2 \pi} d \alpha \int_{0}^{\pi} d \beta \sin \beta \int_{0}^{2 \pi} d \gamma D_{m_{1} j_{1}}^{j_{1}}(\alpha, \beta, \gamma) D_{m_{2}-j_{2}}^{j_{2}}(\alpha, \beta, \gamma) D_{m j_{1}-j_{2}}^{j *}(\alpha, \beta, \gamma), \label{eq:8:1:5}
\end{equation}
where
\[
\clebsch{j_{1}}{j_{1}}{j_{2}}{-j_{2}}{j}{j_{1}-j_{2}}>0 .
\]
Under such a phase definition the Clebsch-Gordan coefficients satisfy the relations
\begin{gather}
\clebsch{j}{m}{0}{0}{j}{m}=1, \quad \clebsch{j_{1}}{j_{1}}{j_{2}}{j_{2}}{j_{1}+j_{2}}{j_{1}+j_{2}}=1, \quad \clebsch{j_{1}}{m_{1}}{j_{2}}{m_{2}}{j_{1}+j_{2}}{m_{1}+m_{2}}>0, \quad \clebsch{j_{1}}{m_{1}}{j_{2}}{-j_{2}}{j}{m}>0, \quad(-1)^{j_{1}+j_{2}-j} \clebsch{j_{1}}{m_{1}}{j_{2}}{j_{2}}{j}{m}>0, \notag\\
(-1)^{j_{1}-m_{1}} \clebsch{j_{1}}{m_{1}}{j_{2}}{m_{2}}{j}{j}>0, \quad(-1)^{j_{2}+m_{2}} \clebsch{j_{1}}{m_{1}}{j_{2}}{m_{2}}{j_{1}-j_{2}}{m}>0 . \label{eq:8:1:6}
\end{gather}
The Clebsch-Gordan coefficients are elements of the unitary matrix which performs direct and inverse transformations between state vectors\index{Clebsch-Gordan coefficient!unitarity}\index{unitarity!of the Clebsch-Gordan coefficients} $\left|j_{1} m_{1} j_{2} m_{2}\right\rangle$ and $\left|j_{1} j_{2} j m\right\rangle$
\begin{equation}
\clebsch{j_{1}}{m_{1}}{j_{2}}{m_{2}}{j}{m}=\braket{j_{1} m_{1} j_{2} m_{2}}{j_{1} j_{2} j m}=\braket{j_{1} j_{2} j m}{j_{1} m_{1} j_{2} m_{2}} . \label{eq:8:1:7}
\end{equation}
The unitarity relation is
\begin{gather}
\sum_{m_{1} m_{2}} \clebsch{j_{1}}{m_{1}}{j_{2}}{m_{2}}{j}{m} \clebsch{j_{1}}{m_{1}}{j_{2}}{m_{2}}{j^{\prime}}{m^{\prime}}=\delta_{j j^{\prime}} \delta_{m m^{\prime}}. \notag\\
\sum_{j(m)} \clebsch{j_{1}}{m_{1}}{j_{2}}{m_{2}}{j}{m} \clebsch{j_{1}}{m_{1}^{\prime}}{j_{2}}{m_{2}^{\prime}}{j}{m}=\delta_{m_{1} m_{1}^{\prime}} \delta_{m_{2} m_{2}^{\prime}}. \label{eq:8:1:8}
\end{gather}
The direct product of two irreducible tensors $\mathfrak{M}_{j_{1} m_{1}}$ and $\mathfrak{N}_{j_{2} m_{2}}$ may be decomposed into irreducible tensors.\index{irreducible tensor product}\isym{Mfrak@$\mathfrak{M}_{J}$, irreducible tensor} The coefficients of this decomposition are just the Clebsch-Gordan coefficients:
\begin{equation}
\mathfrak{M}_{j_{1} m_{1}} \mathfrak{N}_{j_{2} m_{2}}=\sum_{j(m)} \clebsch{j_{1}}{m_{1}}{j_{2}}{m_{2}}{j}{m}\left\{\mathfrak{N}_{j_{1}} \otimes \mathfrak{N}_{j_{2}}\right\}_{j m} . \label{eq:8:1:9}
\end{equation}
The inverse relation is
\begin{equation}
\left\{\mathfrak{M}_{j_{1}} \otimes \mathfrak{N}_{j_{2}}\right\}_{j m}=\sum_{m_{1} m_{2}} \clebsch{j_{1}}{m_{1}}{j_{2}}{m_{2}}{j}{m_{1}} \mathfrak{M}_{j_{1} m_{1}} \mathfrak{R}_{j_{2} m_{2}} . \label{eq:8:1:10}
\end{equation}
\subsection{The Wigner 3jm Symbols}\label{sec:8.1.2}\index{3jm symbol@$3jm$ symbol!definition}
Not infrequently the Wigner $3 j \mathrm{~m}$ symbols \cite{ref110} are used instead of the Clebsch-Gordan coefficients.\isymo{3jm@$\bigl(\begin{smallmatrix} a & b & c\\ \alpha & \beta & \gamma\end{smallmatrix}\bigr)$, Wigner $3jm$ symbol} These symbols possess simpler symmetry properties. The $3 j \mathrm{~m}$ symbols are related to the Clebsch-Gordan coefficients by
\begin{equation}
\threej{j_{1}}{j_{2}}{j_{3}}{m_{1}}{m_{2}}{m_{3}}\label{eq:8:1:11}=(-1)^{j_{3}+m_{3}+2 j_{1}} \frac{1}{\sqrt{2 j_{3}+1}} \clebsch{j_{1}}{-m_{1}}{j_{2}}{-m_{3}}{j_{3}}{m_{3}} .
\end{equation}
The inverse relation is
\begin{equation}
\clebsch{j_{1}}{m_{1}}{j_{2}}{m_{2}}{j_{3}}{m_{3}}=(-1)^{j_{1}-j_{2}+m_{3}} \sqrt{2 j_{3}+1}\threej{j_{1}}{j_{2}}{j_{3}}{m_{1}}{m_{2}}{-m_{3}}\label{eq:8:1:12} .
\end{equation}
The $3 j \mathrm{~m}$ symbol represents the probability amplitude that three angular momenta $\mathrm{j}_{1}, \mathrm{j}_{2}$ and $\mathrm{j}_{3}$ with projections $m_{1}, m_{2}$ and $m_{3}$ are coupled to yield zero angular momentum:
\begin{equation}
\threej{j_{1}}{j_{2}}{j_{3}}{m_{1}}{m_{2}}{m_{3}}\label{eq:8:1:13}=\eta \sum_{j^{\prime} m^{\prime}} \clebsch{j_{1}}{m_{1}}{j_{2}}{m_{3}}{j^{\prime}}{m^{\prime}} \clebsch{j^{\prime}}{m^{\prime}}{j_{3}}{m_{3}}{0}{0}
\end{equation}
The phase factor $\eta=(-1)^{j_{1}-j_{2}+j_{3}}$ is chosen in such a way that any cyclic permutation of columns leaves the $3 j m$ symbol unchanged. Below we shall use, along with the letters $j, m$ other Latin letters ( $a, b, c$, etc.) to denote angular momenta, and Greek letters ( $\alpha, \beta, \gamma$, etc.) to denote momentum projections in the arguments of the $3 j m$ symbols and Clebsch-Gordan coefficients.

\subsection{Regge R-Symbols}\label{sec:8.1.3}\index{Regge R-symbol@Regge $R$-symbol}\index{R-symbol@$R$-symbol|see{Regge $R$-symbol}}
The Wigner $3 j m$ symbol may be represented by a square $3 \times 3$ array\isym{Rik@$R_{i k}$, Regge $R$-symbol}\isym{J-abc@$J=a+b+c$, sum of three momenta} $\left\|R_{i k}\right\|(i, k=1,2,3)$ which is called the Regge $R$-symbol \cite{ref094}
\begin{equation}
\threej{a}{b}{c}{\alpha}{\beta}{\gamma}\label{eq:8:1:14}=\left\|\begin{array}{lll}
R_{11} & R_{12} & R_{13} \\
R_{21} & R_{22} & R_{23} \\
R_{31} & R_{32} & R_{33}
\end{array}\right\|
\end{equation}
where
\begin{equation}
\begin{array}{lll}
R_{11}=-a+b+c, & R_{12}=a-b+c, & R_{13}=a+b-c, \\
R_{21}=a+\alpha, & R_{22}=b+\beta, & R_{23}=c+\gamma, \\
R_{31}=a-\alpha, & R_{32}=b-\beta, & R_{33}=c-\gamma. \label{eq:8:1:15}
\end{array}
\end{equation}
The inverse relations are
\begin{equation}
\begin{array}{ll}
2 a=R_{21}+R_{31}=R_{12}+R_{13}, & 2 \alpha=R_{21}-R_{31}=R_{32}-R_{22}+R_{33}-R_{23}, \\
2 b=R_{22}+R_{32}=R_{11}+R_{13}, & 2 \beta=R_{22}-R_{32}=R_{31}-R_{21}+R_{33}-R_{23}, \\
2 c=R_{23}+R_{33}=R_{11}+R_{12}, & 2 \gamma=R_{23}-R_{33}=R_{31}-R_{21}+R_{32}-R_{22} . \label{eq:8:1:16}
\end{array}
\end{equation}
All nine elements $R_{i k}$ are non-negative integers which satisfy the relations
\begin{equation}
\sum_{i} R_{i k}=J, \quad \sum_{k} R_{i k}=J, \label{eq:8:1:17}
\end{equation}
where
\begin{equation}
J=a+b+c. \label{eq:8:1:18}
\end{equation}
\section{EXPLICIT FORMS OF THE CLEBSCH-GORDAN COEFFICIENTS AND THEIR RELATIONS TO OTHER FUNCTIONS}\label{sec:8.2}\index{Clebsch-Gordan coefficient!explicit forms}
Below we shall assume that the arguments of the Clebsch-Gordan coefficients and $3 j m$ symbols satisfy Eqs.~\ref{eq:8:1:1}-\ref{eq:8:1:3}. In addition we introduce the $\Delta$-symbol defined by\index{Delta-symbol@$\Delta$-symbol}\index{triangle coefficient}\isymg{Delta-abc@$\Delta(a b c)$, triangle coefficient}
\begin{equation}
\Delta(a b c)=\left[\frac{(a+b-c)!(a-b+c)!(-a+b+c)!}{(a+b+c+1)!}\right]^{\frac{1}{2}}. \label{eq:8:2:1}
\end{equation}
The $\Delta$-symbol is invariant under permutations of $a, b, c$. Numerical values of the $\Delta$-symbol are given in Table 8.12 for $\frac{1}{2} \leq a, b, c \leq 5$. If one of the momenta $a, b, c$ equals zero, the $\Delta$-symbol is reduced to
\begin{equation}
\Delta(a a 0)=\frac{1}{\sqrt{2 a+1}}. \label{eq:8:2:2}
\end{equation}
\subsection{Representations of the Clebsch-Gordan Coefficients in the Form of Algebraic Sums}\label{sec:8.2.1}\index{Clebsch-Gordan coefficient!as algebraic sums}
\begin{align}
\clebsch{a}{\alpha}{b}{\beta}{c}{\gamma}= & \delta_{\gamma, \alpha+\beta} \Delta(a b c)[(a+\alpha)!(a-\alpha)!(b+\beta)!(b-\beta)!(c+\gamma)!(c-\gamma)!(2 c+1)]^{\frac{1}{2}}, \notag \\
& \times \sum_{z} \frac{(-1)^{z}}{z!(a+b-c-z)!(a-\alpha-z)!(b+\beta-z)!(c-b+\alpha+z)!(c-a-\beta+z)!}  \label{eq:8:2:3}
\end{align}
(Van der Waerden \cite{ref040}, Racah \cite{ref091}) 
\begin{align}
\clebsch{a}{\alpha}{b}{\beta}{c}{\gamma}= & \delta_{\gamma, \alpha+\beta} \frac{\Delta(a b c)}{(c+a-b)!(c-a+b)!}\left[\frac{(a-\alpha)!(b-\beta)!(c+\gamma)!(c-\gamma)!(2 c+1)}{(a+\alpha)!(b+\beta)!}\right]^{\frac{1}{2}} \notag \\
& \times \sum_{z} \frac{(-1)^{a-\alpha+z}(a+\alpha+z)!(c+b-\alpha-z)!}{z!(a-\alpha-z)!(c-\gamma-z)!(b-c+\alpha+z)!}  \label{eq:8:2:4}
\end{align}
(Racah \cite{ref091}, Fock \cite{ref068})
\begin{align}
\clebsch{a}{\alpha}{b}{\beta}{c}{\gamma}=&\delta_{\gamma, \alpha+\beta} \Delta(a b c)\left[\frac{(c+\gamma)!(c-\gamma)!(2 c+1)}{(a+\alpha)!(a-\alpha)!(b+\beta)!(b-\beta)!}\right]^{\frac{1}{2}}\notag \\ &
 \sum_{z} \frac{(-1)^{b+\beta+z}(c+b+\alpha-z)!(a-\alpha+z)!}{(c-a+b-z)!(c+\gamma-z)!(a-b-\gamma+z)!} \label{eq:8:2:5}
\end{align}
(Wigner \cite{ref043})
\begin{align}
\clebsch{a}{\alpha}{b}{\beta}{c}{\gamma}= & \delta_{\gamma, \alpha+\beta} \frac{\Delta(a b c)}{(a-b+c)!(a+b-c)!}\left[\frac{(a+\alpha)!(a-\alpha)!(b-\beta)!(c+\gamma)!(2 c+1)}{(b+\beta)!(c-\gamma)!}\right]^{\frac{1}{2}}  \notag\\
& \times \sum_{z} \frac{(-1)^{b+\beta+x}(2 c-z)!(a+b-c+z)!}{z!(c-a+b-z)!(c+\gamma-z)!(a-c-\beta+z)!}. \label{eq:8:2:6}
\end{align}
(Majumdar \cite{ref082})
\begin{align}
\clebsch{a}{\alpha}{b}{\beta}{c}{\gamma}= & \delta_{\gamma, \alpha+\beta} \frac{\Delta(a b c)(a+b+c+1)!}{(a-b+c)!}\left[\frac{(a-\alpha)!(c+\gamma)!(2 c+1)}{(a+\alpha)!(b+\beta)!(b-\beta)!(c-\gamma)!}\right]^{\frac{1}{2}}. \notag\\
& \times \sum_{z} \frac{(-1)^{b+\beta+z}(2 c-z)!(c+b+\alpha-z)!}{z!(c-a+b-z)!(c+\gamma-z)!(a+b+c+1-z)!}. \label{eq:8:2:7}
\end{align}
(Bandzaitis and Yutsis \cite{ref051})
\begin{align}
\clebsch{a}{\alpha}{b}{\beta}{c}{\gamma}= & \frac{\delta_{\gamma, \alpha+\beta}}{\Delta(a b c)}\left[\frac{(a+\alpha)!(a-\alpha)!(c+\gamma)!(c-\gamma)!(2 c+1)}{(b+\beta)!(b-\beta)!}\right]^{\frac{1}{2}}  \notag\\
& \times \sum_{z} \frac{(-1)^{a-\alpha+z}(a+b-\gamma-z)!(b+c-\alpha-z)!}{z!(a-\alpha-z)!(c-\gamma-z)!(a+b+c+1-z)!}. \label{eq:8:2:8}
\end{align}
(Bandzaitis and Yutsis \cite{ref051}).\\
In Eqs.~\eqref{eq:8:2:3}-\eqref{eq:8:2:8} the summation index $z$ assumes integer values for which all the factorial arguments are nonnegative.

\subsection{Quasi-Binomial Representation of the Clebsch-Gordan Coefficients}\label{sec:8.2.2}\index{Clebsch-Gordan coefficient!quasi-binomial representation}\index{quasi-binomial}
The Clebsch-Gordan coefficients can be presented in the form of quasi-binomials $(u \pm v)^{(k)}$ \cite{ref048}
\begin{equation}
\left.\left.\clebsch{a}{\alpha}{b}{\beta}{c}{\gamma}=\delta_{\gamma, \alpha+\beta} \frac{\Delta(a b c)}{(a+b-c)!}\left[\frac{(c-\gamma)!(c+\gamma)!(2 c+1)}{(a-\alpha)!(a+\alpha)!(b-\beta)!(b+\beta)!}\right]^{\frac{1}{2}} \right\rvert\,(a+\alpha)(b-\beta)-(a-\alpha)(b+\beta)\right]^{(a+b-c)} . \label{eq:8:2:9}
\end{equation}
For calculating the quasi-binomial one should use the binomial formula with all powers replaced by quasipowers.
\begin{equation}
(u \pm v)^{(k)}=\sum_{z}( \pm 1)^{z}\binom{k}{z} u^{(k-z)} v^{(z)}. \label{eq:8:2:10}
\end{equation}
Quasi-powers (or generalized powers) are defined by\index{quasi-power}\index{generalized power|see{quasi-power}}\isym{quasipower@$u^{(z)}$, quasi-power}
\begin{gather}
u^{(z)} \equiv u^{(1)(z)}=\frac{\Gamma(u+1)}{\Gamma(u-z+1)}, \notag\\
u^{(-1)(z)}=(u+z)^{(z)}=\frac{\Gamma(u+z+1)}{\Gamma(u+1)}. \label{eq:8:2:11}
\end{gather}
If $z$ is integer and positive, we get
\begin{align}
u^{(z)} & =u(u-1) \ldots(u-z+1), \notag\\
u^{(-1)(z)} & =(u+1)(u+2) \ldots(u+z). \label{eq:8:2:12}
\end{align}
\subsection{Clebsch-Gordan Coefficients and Finite Differences}\label{sec:8.2.3}\index{Clebsch-Gordan coefficient!and finite differences}\index{finite differences}
The Clebsch-Gordan coefficient are related to finite differences \cite{ref041}
\begin{gather}
\clebsch{a}{\alpha}{b}{\beta}{c}{\gamma}=\delta_{\gamma, \alpha+\beta}(-1)^{a-c+\beta}\left[\frac{(c+\gamma)!(a+b-c)!(2 c+1)}{(c-\gamma)!(c+a-b)!(c-a+b)!(a+b+c+1)!}\right]^{\frac{1}{2}}  \notag\\
\times\left[\frac{(a-\alpha)^{(a-b-\gamma)}}{(a+\alpha)^{(a-b+\gamma)}}\right]^{\frac{1}{2}} \Delta_{\alpha}^{c-\gamma}\left[\frac{(a+\alpha)^{(c+a-b)}}{(a-\alpha)^{(a-b-c)}}\right]. \label{eq:8:2:13}
\end{gather}
where the difference of order $k$ with respect to the argument $a$ is given by\isymg{Delta-fd@$\Delta_{a}^{k}$, finite-difference operator}
\begin{equation}
\Delta_{a}^{k} f(a)=\sum_{n=0}^{k}(-1)^{k-n} \frac{k!}{n!(k-n)!} f(a+n). \label{eq:8:2:14}
\end{equation}
In particular,
\begin{gather}
\Delta_{a} f(a)=f(a+1)-f(a), \notag\\
\Delta_{a} a^{(n)}=n a^{(n-1)} . \label{eq:8:2:15}
\end{gather}
Finite differences are analogous to differentials, while quasi-powers are analogous to ordinary powers. From this it can be concluded that the Clebsch-Gordan coefficients are analogous to the Wigner $D$-functions [20, 58].\index{Wigner D-function@Wigner $D$-function}\isym{DJMM@$D_{M M^{\prime}}^{J}$, Wigner $D$-function} To illustrate this similarity let us introduce the notations
\begin{equation}
j=b, \quad \mu=c-a, \quad \nu=\beta, \quad P=\frac{1}{2}(a+b+c), \quad Q=a+\alpha-\frac{1}{2}(a+b+c) . \label{eq:8:2:16}
\end{equation}
\footnotetext{${ }^{2}$ This follows from the fact that the Clebsch-Gordan coefficients may be expressed in terms of orthogonal polynomials of a discrete variable (Hahn polynomials); see Refs.~\cite{ref142, ref143}.
}Now, Eq.~\eqref{eq:8:2:3} may be rewritten in terms of quasi-powers
\begin{align}
\sqrt{\frac{a+b+c+1}{2 c+1}} \clebsch{a}{\alpha}{b}{\beta}{c}{\gamma}= & {\left[\frac{(j+\mu)!(j-\mu)!(j+\nu)!(j-\nu)!}{(2 P)^{(2 j)}(P-Q)^{(j+\mu)}(P-Q)^{(j+\nu)}(P+Q)^{(-\mu-\nu)}}\right]^{\frac{1}{2}} }  \notag\\
& \times \sum_{x}(-1)^{z} \frac{(P-Q)^{(j+\mu+z)}(P+Q)^{(j-\mu-z)}}{z!(j-\mu-z)!(j+\nu-z)!(\mu-\nu+z)!}. \label{eq:8:2:17}
\end{align}
Replacing quasi-powers in \eqref{eq:8:2:17} by ordinary powers and introducing the angle $\vartheta$ according to
\begin{equation}
\cos \theta=\frac{Q}{P}, \quad \cos \frac{\theta}{2}=\sqrt{\frac{P+Q}{2 P}}, \quad \sin \frac{\theta}{2}=\sqrt{\frac{P-Q}{2 P}}, \label{eq:8:2:18}
\end{equation}
we see that \eqref{eq:8:2:17} reduces to the following relation for the Wigner $D$-function
\begin{align}
& d_{\nu \mu}^{j}(\theta)=\left[\frac{(j+\mu)!(j-\mu)!(j+\nu)!(j-\nu)!}{(2 P)^{2 j}(P-Q)^{j+\mu}(P-Q)^{j+\nu}(P+Q)^{-\mu-\nu}}\right]^{\frac{1}{2}}  \notag\\
& \times \sum_{x}(-1)^{x} \frac{(P-Q)^{j+\mu+x}(P+Q)^{j-\mu-x}}{z!(j-\mu-z)!(j+\nu-z)!(\mu-\nu+z)!}. \label{eq:8:2:19}
\end{align}
\eqref{eq:8:2:19} is similar to Eq.~\eqref{eq:4:3:4}. The analogy between \eqref{eq:8:2:17} and \eqref{eq:8:2:19} permits us to connect the recursion relations, symmetry properties, etc. of the Wigner D-functions and the Clebsch-Gordan coefficients \cite{ref058}. Since for $u \gg z, u \gg 1$ quasi-powers became asymptotically equal to ordinary powers, a Clebsch-Gordan coefficient turns into a $D$-function in the large-momentum limit (\eqref{eq:8:9:1}).

\subsection{Expressions for the Clebsch-Gordan Coefficients in Terms of the Binomial Coefficients \cite{ref126}}\label{sec:8.2.4}\index{Clebsch-Gordan coefficient!in terms of binomial coefficients}\index{binomial coefficients}
\begin{equation}
\clebsch{a}{\alpha}{b}{\beta}{c}{\gamma}=\delta_{\gamma, \alpha+\beta}\left[\frac{\binom{2 a}{J-2 c}\binom{2 b}{J-2 c}}{\binom{J+1}{J-2 c}\binom{2 a}{a-\alpha}\binom{2 b}{b-\beta}\binom{2 c}{c-\gamma}}\right]^{\frac{1}{2}} \sum_{x}(-1)^{x}\binom{J-2 c}{z}\binom{J-2 b}{a-\alpha-z}\binom{J-2 a}{b+\beta-z}, \label{eq:8:2:20}
\end{equation}
where $J=a+b+c$.

\subsection{Representations of the Clebsch-Gordan Coefficients in Terms of of the Hypergeometric Functions}\label{sec:8.2.5}\index{Clebsch-Gordan coefficient!and hypergeometric functions}\index{hypergeometric function!generalized}
The Clebsch-Gordan coefficients may be expressed via the generalized hypergeometric functions $\Fhg{3}{2}$ of unit argument,\isym{F32@$\Fhg{3}{2}$, generalized hypergeometric function}
\begin{gather}
\clebsch{a}{\alpha}{b}{\beta}{c}{\gamma}=\delta_{\gamma, \alpha+\beta} \frac{\Delta(a b c)}{(a+b-c)!(-b+c+\alpha)!(-a+c-\beta)!}\left[\frac{(a+\alpha)!(b-\beta)!(c+\gamma)!(c-\gamma)!(2 c+1)}{(a-\alpha)!(b+\beta)!}\right]^{\frac{1}{2}}  \notag\\
\times\Fhg{3}{2}\left[\left.\begin{array}{c}
-a-b+c,-a+\alpha,-b-\beta \\
-a+c-\beta+1,-b+c+\alpha+1
\end{array} \right\rvert\, 1\right], \label{eq:8:2:21}
\end{gather}
\begin{align}
& \begin{aligned}
\clebsch{a}{\alpha}{b}{\beta}{c}{\gamma}= & \delta_{\gamma, \alpha+\beta}(-1)^{a-\alpha} \frac{\Delta(a b c)(b+c-\alpha)!}{(a-b+c)!(-a+b+c)!(b-c+\alpha)!}\left[\frac{(a+\alpha)!(b-\beta)!(c+\gamma)!(2 c+1)}{(a-\alpha)!(b+\beta)!(c-\gamma)!}\right]^{\frac{1}{2}} \\
& \times\Fhg{3}{2}\left[\left.\begin{array}{c}
a+\alpha+1,-a+\alpha,-c+\gamma \\
-b-c+\alpha, b-c+\alpha+1
\end{array} \right\rvert\, 1\right],
\end{aligned}, \\
& \begin{aligned}
\clebsch{a}{\alpha}{b}{\beta}{c}{\gamma}= & \delta_{\gamma, \alpha+\beta}(-1)^{b+\beta} \frac{\Delta(a b c)(b+c+\alpha)!}{(-a+b+c)!(a-b-\gamma)!}\left[\frac{(a-\alpha)!(c-\gamma)!(2 c+1)}{(a+\alpha)!(b+\beta)!(b-\beta)!(c+\gamma)!}\right]^{\frac{1}{2}} \\
& \quad \times\Fhg{3}{2}\left[\left.\begin{array}{c}
a-b-c, a-\alpha+1,-c-\gamma \\
a-b-\gamma+1,-b-c-\alpha \\
\end{array} \right\rvert\, 1\right]\\
&\qquad 
(\text { Rose \cite{ref030}})
\end{aligned}, \label{eq:8:2:23}\\
& \clebsch{a}{\alpha}{b}{\beta}{c}{\gamma}=\delta_{\gamma, \alpha+\beta}(-1)^{b+\beta} \frac{\Delta(a b c)(2 c)!}{(a-b+c)!(-a+b+c)!(a-c-\beta)!}\left[\frac{(a+\alpha)!(a-\alpha)!(b-\beta)!(2 c+1)}{(b+\beta)!(c+\gamma)!(c-\gamma)!}\right]^{\frac{1}{2}}  \notag\\
&  \label{eq:8:2:24}\\
& \quad \times\Fhg{3}{2}\left[\begin{array}{c}
a-b-c, a+b-c+1,-c-\gamma \mid 1], \\
-2 c, a-c-\beta+1
\end{array}\right. \notag\\
& \clebsch{a}{\alpha}{b}{\beta}{c}{\gamma}=\delta_{\gamma, \alpha+\beta}(-1)^{b+\beta} \frac{\Delta(a b c)(2 c)!(b+c+\alpha)!}{(a-b+c)!(-a+b+c)!}\left[\frac{(a+\alpha)!(b+\beta)!(b-\beta)!(c+\gamma)!(c-\gamma)!}{(a-\alpha)}\right]^{\frac{1}{2}}  \notag\\
& \quad \times\Fhg{3}{2}\left[\begin{array}{c}
a-b-c,-a-b-c-1,-c-\gamma \mid 1], \\
-2 c,-b-c-\alpha
\end{array}\right. \label{eq:8:2:25}\\
& \clebsch{a}{\alpha}{b}{\beta}{c}{\gamma}=\delta_{\gamma, \alpha+\beta}(-1)^{a-\alpha} \frac{(a+b-\gamma)!(b+c-\alpha)!}{\Delta(a b c)(a+b+c+1)!}\left[\frac{(a+\alpha)!(c+\gamma)!(2 c+1)}{(a-\alpha)!(b+\beta)!(b-\beta)!(c-\gamma)!}\right]^{\frac{1}{2}}  \notag\\
& \quad \times\Fhg{3}{2}\left[\begin{array}{c}
-a-b-c-1,-a+\alpha,-c+\gamma \mid 1 \\
-a-b+\gamma,-b-c+\alpha
\end{array}\right] . \label{eq:8:2:26}
\end{align}
Equations \eqref{eq:8:2:21}-\eqref{eq:8:2:26} may be associated with Eqs.~\eqref{eq:8:2:3}-\eqref{eq:8:2:8}, respectively.\\[0pt]
According to Ref.~\cite{ref036}, all $\Fhg{3}{2}$-functions of unit argument, which are given by finite sums, may be expressed in terms of the Clebsch-Gordan coefficients. Any Clebsch-Gordan coefficient may also be represented in the form of derivatives of the hypergeometric functions \cite{ref046}:
\begin{gather}
\clebsch{a}{\alpha}{b}{\beta}{c}{\gamma}=\delta_{\gamma, \alpha+\beta}(-1)^{a+b-c} \frac{\Delta(a b c)}{(a+b-c)!(-a+b+c)!(a-b+\gamma)!}\left[\frac{(a+\alpha)!(a-\alpha)!(b+\beta)!(c+\gamma)!(2 c+1)}{(b-\beta)!(c-\gamma)!}\right]^{\frac{1}{2}}  \notag\\
\times\left\{\frac{d^{b-\beta}}{(d t)^{b-\beta}}\left[(1-t)^{a+b-c} F(a-b-c,-c+\gamma ; a-b+\gamma+1 ; t)\right]\right\}_{t=0}. \label{eq:8:2:27}
\end{gather}
\subsection{Representations of the $3 j \mathrm{~m}$ Symbols in the Form of Algebraic Sums}\label{sec:8.2.6}\index{3jm symbol@$3jm$ symbol!as algebraic sums}
The explicit form of the Wigner $3 j m$ symbols may be obtained from Eqs.~\eqref{eq:8:2:3}-\eqref{eq:8:2:8} for the Clebsch-Gordan coefficients. Rewriting these equations in terms of the $R$-symbol elements, we get the following relations for\\
the $3 j m$ symbols \cite{ref045}
\begin{align}
 \threej{a}{b}{c}{\alpha}{\beta}{\gamma}&=(-1)^{R_{21}+R_{32}}\left[\frac{\prod_{i, j=1}^{3}\left(R_{i j}\right)!}{(J+1)!}\right]^{\frac{1}{2}} \notag \\
& \times \sum_{x}(-1)^{x} \frac{1}{z!\left(R_{13}-z\right)!\left(R_{22}-z\right)!\left(R_{31}-z\right)!\left(R_{33}-R_{22}+z\right)!\left(R_{11}-R_{22}+z\right)!}, \label{eq:8:2:28}\\
\threej{a}{b}{c}{\alpha}{\beta}{\gamma}&=(-1)^{R_{21}+R_{31}+R_{32}}\left[\frac{R_{13}!R_{23}!R_{31}!R_{32}!R_{33}!}{(J+1)!R_{11}!R_{12}!R_{21}!R_{22}!}\right]^{\frac{1}{2}}  \notag\\
 &\times \sum_{x}(-1)^{x} \frac{\left(R_{21}+z\right)!\left(R_{11}+R_{31}-z\right)!}{z!\left(R_{31}-z\right)!\left(R_{23}-z\right)!\left(R_{13}-R_{31}+z\right)!}, \label{eq:8:2:29}\\
 \threej{a}{b}{c}{\alpha}{\beta}{\gamma}&=(-1)^{R_{21}+R_{22}+R_{32}}\left[\frac{R_{11}!R_{12}!R_{13}!R_{23}!R_{33}!}{(J+1)!R_{21}!R_{22}!R_{31}!R_{32}!}\right]^{\frac{1}{2}} \sum_{z}(-1)^{z} \frac{\left(R_{31}+z\right)!\left(R_{32}+R_{33}-z\right)!}{z!\left(R_{11}-z\right)!\left(R_{33}-z\right)!\left(R_{12}-R_{33}+z\right)!}, \\
 \threej{a}{b}{c}{\alpha}{\beta}{\gamma}&=(-1)^{R_{21}+R_{22}+R_{32}}\left[\frac{R_{11}!R_{21}!R_{31}!R_{32}!R_{33}!}{(J+1)!R_{12}!R_{13}!R_{22}!R_{33}!}\right]^{\frac{1}{2}} \sum_{z}(-1)^{z} \frac{\left(R_{13}+z\right)!\left(R_{23}+R_{33}-z\right)!}{z!\left(R_{11}-z\right)!\left(R_{33}-z\right)!\left(R_{21}-R_{33}+z\right)!}, \label{eq:8:2:31}\\
 \threej{a}{b}{c}{\alpha}{\beta}{\gamma}&=(-1)^{R_{21}+R_{22}+R_{32}}\left[\frac{(J+1)!R_{11}!R_{13}!R_{31}!R_{39}!}{R_{12}!R_{21}!R_{22}!R_{23}!R_{32}!}\right]^{\frac{1}{2}} \sum_{z}(-1)^{z} \frac{\left(R_{11}+R_{21}-z\right)!\left(R_{23}+R_{33}-z\right)!}{z!\left(R_{11}-z\right)!\left(R_{33}-z\right)!(J+1-z)!}, \label{eq:8:2:32}\\
 \threej{a}{b}{c}{\alpha}{\beta}{\gamma}&=(-1)^{R_{21}+R_{31}+R_{32}}\left[\frac{(J+1)!R_{21}!R_{23}!R_{31}!R_{39}!}{R_{11}!R_{12}!R_{13}!R_{22}!R_{32}!}\right]^{\frac{1}{2}} \sum_{z}(-1)^{z} \frac{\left(R_{11}+R_{31}-z\right)!\left(R_{13}+R_{23}-z\right)!}{z!\left(R_{23}-z\right)!\left(R_{31}-z\right)!(J+1-z)!}, \label{eq:8:2:33}
\end{align}
where $J=a+b+c$.

\subsection{Quasi-binomial Representations of the $3jm$ Symbols}\label{sec:8.2.7}\index{3jm symbol@$3jm$ symbol!quasi-binomial representations}
\begin{equation}
\threej{a}{b}{c}{\alpha}{\beta}{\gamma}\label{eq:8:2:34}=(-1)^{R_{21}+R_{32}}\left[\frac{R_{21}^{(-1)\left(R_{31}-R_{13}\right)} R_{31}^{(-1)\left(R_{32}-R_{13}\right)} R_{32}^{(-1)\left(R_{22}-R_{13}\right)} R_{22}^{(-1)\left(R_{21}-R_{13}\right)}}{(J+1)!R_{13}!}\right]^{\frac{1}{4}}(u-v)^{\left(R_{13}\right)} .
\end{equation}
The quasi-binomial may be evaluated according to Eqs.~\eqref{eq:8:2:10} and \eqref{eq:8:2:11}. Variables $u$ and $v$ may be chosen in different ways \cite{ref045}:
\begin{align}
(a)& \qquad u=R_{21}^{(1)} R_{32}^{(1)}, && v=R_{22}^{(1)} R_{31}^{(1)},\notag\\
(b)& \qquad u=R_{21}^{(1)} R_{11}^{(-1)}, && v=R_{22}^{(1)} R_{12}^{(-1)},\notag\\
(c)& \qquad u=R_{11}^{(-1)} R_{23}^{(-1)}, && v=R_{32}^{(1)}(J+1)^{(1)}.
\end{align}

\section{INTEGRAL REPRESENTATIONS}\label{sec:8.3}\index{Clebsch-Gordan coefficient!integral representations}\index{integral representations!of the Clebsch-Gordan coefficients}
\subsection{Integrals Involving Algebraic Functions}\label{sec:8.3.1}\index{integral representations!with algebraic functions}
The Clebsch-Gordan coefficients may be represented by the following integrals \cite{ref041} ( $J=a+b+c$ )
\begin{gather}
\clebsch{a}{\alpha}{b}{\beta}{c}{\gamma}=\frac{(-1)^{a-c+\beta}}{2^{J+1}}\left[\frac{(c+\gamma)!(J-2 c)!(J+1)!(2 c+1)}{(a-\alpha)!(a+\alpha)!(b-\beta)!(b+\beta)!(c-\gamma)!(J-2 a)!(J-2 b)!}\right]^{\frac{1}{2}}  \notag\\
\times \int_{-1}^{1}(1-x)^{a-\alpha}(1+x)^{b-\beta} \frac{d^{c-\gamma}}{(d x)^{c-\gamma}}\left[(1-x)^{J-2 a}(1+x)^{J-2 b}\right] d x, \label{eq:8:3:1}\\
\clebsch{a}{\alpha}{b}{\beta}{c}{\gamma}=\frac{(-1)^{a-c+\beta}}{2^{J+1}}\left[\frac{(J-2 b)!(J-2 c)!(J+1)!(2 c+1)}{(a-\alpha)!(a+\alpha)!(b-\beta)!(b+\beta)!(c-\gamma)!(c+\gamma)!(J-2 a)!}\right]^{\frac{1}{2}}  \notag\\
\times \int_{-1}^{1}(1-x)^{b+\beta}(1+x)^{b-\beta} \frac{d^{J-2 a}}{(d x)^{J-2 a}}\left[(1-x)^{c-\gamma}(1+x)^{c+\gamma}\right] d x. \label{eq:8:3:2}
\end{gather}
In particular, if $\alpha=\beta=\gamma=0$, then
\begin{equation}
\clebsch{a}{0}{b}{0}{c}{0}=\frac{(-1)^{a-c}}{2^{J+1} a!b!c!}\left[\frac{(J-2 b)!(J-2 c)!(J+1)!(2 c+1)}{(J-2 a)!}\right]^{\frac{1}{2}} \int_{-1}^{1}\left(1-x^{2}\right)^{b} \frac{d^{J-2 a}}{(d x)^{J-2 a}}\left[\left(1-x^{2}\right)^{c}\right] d x. \label{eq:8:3:3}
\end{equation}
\subsection{Integrals Involving the Wigner D-Functions}\label{sec:8.3.2}\index{integral representations!with the Wigner $D$-functions}\isym{dR@$d R$, invariant volume element of the rotation group}
\begin{equation}
\clebsch{a}{\alpha}{b}{\beta}{c}{\gamma}=\frac{1}{8 \pi^{2}}\left[\frac{(a+b+c+1)!(a+b-c)!(2 c+1)}{(2 a)!(2 b)!}\right]^{\frac{1}{2}} \int d R D_{\alpha a}^{a}(R) D_{\beta-b}^{b}(R) D_{\gamma a-b}^{c *}(R), \label{eq:8:3:4}
\end{equation}
where $R$ means the set of the Euler angles (see Sec.~\ref{sec:4.12})
\begin{equation}
\begin{aligned}
\clebsch{a}{\alpha}{b}{\beta}{c}{\gamma}= & \frac{(-1)^{a+c-\alpha-\gamma}}{2^{a+b+1}}\left[\frac{(J-2 c)!(J+1)!(2 c+1)}{(a+\alpha)!(a-\alpha)!(b+\beta)!(b-\beta)!}\right]^{\frac{1}{2}} \\
& \times \int_{-1}^{1}(\sin \theta)^{a+b}\left(\frac{1-\cos \theta}{1+\cos \theta}\right)^{\frac{\alpha-\beta}{2}} d_{b-a \gamma}^{c}(\theta) d(\cos \theta)\\
&\qquad 
\text{(Akim and Levin, \cite{ref046})}
\end{aligned}
\end{equation}

\subsection{Integral Involving the Spherical Harmonics}\label{sec:8.3.3}\index{integral representations!with the spherical harmonics}\index{spherical harmonics}\isym{Ylm@$Y_{l m}(\vartheta, \varphi)$, spherical harmonic}
If $a, b, c$ are integers and $J=a+b+c=2 g$ is even (i.e., $g$ is integer), then
\begin{equation}
\clebsch{a}{\alpha}{b}{\beta}{c}{\gamma}=(-1)^{g-c}\left[\frac{4 \pi}{(2 a+1)(2 b+1)}\right]^{\frac{1}{2}} \frac{(g-a)!(g-b)!(g-c)!}{\Delta(a b c) g!} \int_{0}^{2 \pi} d \varphi \int_{0}^{\pi} d \vartheta \sin \vartheta Y_{a \alpha}(\vartheta, \varphi) Y_{b \beta}(\vartheta, \varphi) Y_{c \gamma}^{*}(\vartheta, \varphi). \label{eq:8:3:6}
\end{equation}
\subsection{Integral Representations for Products of the Clebsch-Gordan Coefficients}\label{sec:8.3.4}\index{integral representations!for products of Clebsch-Gordan coefficients}\index{Legendre polynomials}\isym{Pl@$P_{l}(x)$, Legendre polynomial}
\begin{equation}
\clebsch{a}{\alpha}{b}{\beta}{c}{\gamma} \clebsch{a}{\alpha^{\prime}}{b}{\beta^{\prime}}{c}{\gamma^{\prime}}=\frac{2 c+1}{8 \pi^{2}} \int d R D_{\alpha \alpha^{\prime}}^{a}(R) D_{\beta \beta^{\prime}}^{b}(R) D_{\gamma \gamma^{\prime}}^{c *}(R). \label{eq:8:3:7}
\end{equation}
In particular,
\begin{equation}
\left[\clebsch{a}{0}{b}{0}{c}{0}\right]^{2}=\frac{2 c+1}{2} \int_{0}^{\pi} d \theta \sin \vartheta P_{a}(\cos \vartheta) P_{b}(\cos \vartheta) P_{c}(\cos \vartheta). \label{eq:8:3:8}
\end{equation}
\section{SYMMETRY PROPERTIES}\label{sec:8.4}\index{vector-addition coefficients!symmetry properties}\index{symmetry!of the vector-addition coefficients}
The simplest way to formulate the symmetry properties of the Clebsch-Gordan coefficients and 3jm symbols is to use the representations of these coefficients in the form of the Regge $R$-symbols \cite{ref094} (see Sec.~\ref{sec:8.1.3}).

\subsection{Symmetry Properties of the $R$-Symbols}\label{sec:8.4.1}\index{Regge R-symbol@Regge $R$-symbol!symmetry properties}
The $R$-symbol possesses the following symmetry properties \cite{ref094}.\\
\subsubsection{Permutations of columns}
\begin{equation}
\left\|\begin{array}{lll}
R_{11} & R_{12} & R_{13}\\
R_{21} & R_{22} & R_{23} \\
R_{31} & R_{32} & R_{33}
\end{array}\right\|=\varepsilon\left\|\begin{array}{lll}
R_{1 i} & R_{1 k} & R_{1 i} \\
R_{2 i} & R_{2 k} & R_{2 l} \\
R_{3 i} & R_{3 k} & R_{3 l}
\end{array}\right\|,   \label{eq:8:4:1}
\end{equation}
where
\begin{equation}
c= \begin{cases}+1 & \text { for cyclic permutation (even number of permutations) }  \\ (-1)^{J} & \text { for non-cyclic permutation (odd number of permutations) }\end{cases}, \label{eq:8:4:2}
\end{equation}
and
\[
J \equiv \sum_{k} R_{i k}=\sum_{i} R_{i k} ;
\]
\subsubsection{Permutations of rows}
\begin{equation}
\left\|\begin{array}{lll}
R_{11} & R_{12} & R_{13}≈\\
R_{21} & R_{22} & R_{23} \\
R_{31} & R_{32} & R_{33}
\end{array}\right\|=\left\|\begin{array}{lll}
R_{i 1} & R_{i 2} & R_{i 3} \\
R_{k 1} & R_{k 2} & R_{k 3} \\
R_{i 1} & R_{i 2} & R_{i 3}
\end{array}\right\|, 
\end{equation}
\subsubsection{Transposition}
\begin{equation}
\left\|\begin{array}{lll}
R_{11} & R_{12} & R_{13} \\
R_{21} & R_{22} & R_{23} \\
R_{31} & R_{32} & R_{33}
\end{array}\right\|=\left\|\begin{array}{lll}
R_{11} & R_{21} & R_{31} \\
R_{12} & R_{22} & R_{32} \\
R_{13} & R_{23} & R_{33}
\end{array}\right\|.  \label{eq:8:4:4}
\end{equation}
These symmetry transformations relate $6 \times 6 \times 2=72$ generally different coefficients.

\subsection{Symmetry Properties of the $3 j m$ Symbols}\label{sec:8.4.2}\index{3jm symbol@$3jm$ symbol!symmetry properties}
The above relations are equivalent to the following symmetry properties of the $3 j m$ symbol.

\subsubsection{Classical symmetries \cite{ref110}}
\begin{enumerate}[label=(\roman*)]
\item Permutations of columns (corresponding to permutations of columns of the $R$-symbol):
\begin{align}
\threej{a}{b}{c}{\alpha}{\beta}{\gamma}=\threej{b}{c}{a}{\beta}{\gamma}{\alpha}=\threej{c}{a}{b}{\gamma}{\alpha}{\beta} & =(-1)^{a+b+c}\threej{a}{c}{b}{\alpha}{\gamma}{\beta}  \notag\\
& =(-1)^{\alpha+\beta+\gamma}\threej{b}{a}{c}{\beta}{\alpha}{\gamma}=(-1)^{\alpha+b+c}\threej{c}{b}{a}{\gamma}{\beta}{\alpha}. \label{eq:8:4:5}
\end{align}
\item Change of signs of momentum projections (corresponding to the permutation of the second and third rows of the $R$-symbol)
\begin{equation}
\threej{a}{b}{c}{\alpha}{\beta}{\gamma}\label{eq:8:4:6}=(-1)^{a+b+c}\threej{a}{b}{c}{-\alpha}{-\beta}{-\gamma}
\end{equation}
\subsubsection{Regge symmetries \cite{ref094}}
\item Replacement of arguments:
\begin{equation}
\threej{a}{b}{c}{\alpha}{\beta}{\gamma}\label{eq:8:4:7}=(-1)^{a+b+c}\threej{\frac{b+c+\alpha}{2}}{\frac{a+c+\beta}{2}}{\frac{a+b+\gamma}{2}}{a-\frac{b+c-\alpha}{2}}{b-\frac{a+c-\beta}{2}}{c-\frac{a+b-\gamma}{2}}
\end{equation}
(corresponding to the permutation of the first and third rows of the $R$-symbol) and\\
\item
\begin{equation}
\threej{a}{b}{c}{\alpha}{\beta}{\gamma}\label{eq:8:4:8}=\threej{a}{\frac{b+c-\alpha}{2}}{\frac{b+c+\alpha}{2}}{-b+c}{\frac{b-c-\alpha}{2}-\gamma}{\frac{b-c+\alpha}{2}+\gamma}
\end{equation}
This corresponds to the transposition of the $R$-symbol. According to the above symmetry properties, 72 formally different $3 j m$ symbols have the same absolute values. One can decompose these symbols into six groups ( 12 coefficients in each group). The $3 j m$ symbols which belong to different groups correspond to triangles of different forms but the same perimeter $J=a+b+c$. On the other hand, all $3 j m$ symbols which belong to one group are related by the classical symmetry properties \eqref{eq:8:4:5} and \eqref{eq:8:4:6}. In \eqref{eq:8:4:9} we list one $3 j m$ symbol from each group.
\begin{align}
& \threej{a}{b}{c}{\alpha}{\beta}{\gamma}=\threej{a}{\frac{b+c-\alpha}{2}}{\frac{b+c+\alpha}{2}}{c-b}{\frac{b+\beta-c-\gamma}{2}}{\frac{b-\beta-c+\gamma}{2}}=\threej{b}{\frac{a+c-\beta}{2}}{\frac{a+c+\beta}{2}}{a-c}{\frac{-a-\alpha+c+\gamma}{2}}{\frac{-a+\alpha+c-\gamma}{2}}  \notag\\
& =\threej{c}{\frac{a+b-\gamma}{2}}{\frac{a+b+\gamma}{2}}{b-a}{\frac{a+\alpha-b-\beta}{2}}{\frac{a-\alpha-b+\beta}{2}}=\threej{\frac{b+c-\alpha}{2}}{\frac{a+c-\beta}{2}}{\frac{a+b-\gamma}{2}}{a-\alpha-\frac{b+c-\alpha}{2}}{b-\beta-\frac{a+c-\beta}{2}}{c-\gamma-\frac{a+b-\gamma}{2}}  \notag\\
& =\threej{\frac{b+c+\alpha}{2}}{\frac{a+c+\beta}{2}}{\frac{a+b+\gamma}{2}}%
{\frac{b+c+\alpha}{2}-a-\alpha}{\frac{a+c+\beta}{2}-b-\beta}{\frac{a+b+\gamma}{2}-c-\gamma}
 . \label{eq:8:4:9}
\end{align}
\end{enumerate}
\subsection{Symmetry Properties of the Clebsch-Gordan Coefficients}\label{sec:8.4.3}\index{Clebsch-Gordan coefficient!symmetry properties}
\begin{enumerate}[label=(\roman*)]
\item
  \begin{align}
\clebsch{a}{\alpha}{b}{\beta}{c}{\gamma}=(-1)^{a+b-c} \clebsch{b}{\beta}{a}{\alpha}{c}{\gamma} & =(-1)^{a-\alpha} \sqrt{\frac{2 c+1}{2 b+1}} \clebsch{a}{\alpha}{c}{-\gamma}{b}{-\beta}=(-1)^{a-\alpha} \sqrt{\frac{2 c+1}{2 b+1}} \clebsch{c}{\gamma}{a}{-\alpha}{b}{\beta}  \notag\\
& =(-1)^{b+\beta} \sqrt{\frac{2 c+1}{2 a+1}} \clebsch{c}{-\gamma}{b}{\beta}{a}{-\alpha}=(-1)^{b+\beta} \sqrt{\frac{2 c+1}{2 a+1}} \clebsch{b}{-\beta}{c}{\gamma}{a}{\alpha}. \label{eq:8:4:10}
\end{align}
\item
\begin{equation}
\clebsch{a}{\alpha}{b}{\beta}{c}{\gamma}=(-1)^{a+b-c} \clebsch{a}{-\alpha}{b}{-\beta}{c}{-\gamma}. \label{eq:8:4:11}
\end{equation}
\item
\begin{equation}
\clebsch{a}{\alpha}{b}{\beta}{c}{\gamma}=\clebsch{a^{\prime}}{\alpha^{\prime}}{b^{\prime}}{\beta^{\prime}}{c^{\prime}}{\gamma^{\prime}}, \label{eq:8:4:12}
\end{equation}
where
\begin{align}
2 a^{\prime} & =(a+\alpha)+(b+\beta), & 2 a & =\left(a^{\prime}+a^{\prime}\right)+\left(b^{\prime}+\beta^{\prime}\right), \notag\\
2 \alpha^{\prime} & =(a+\alpha)-(b+\beta), & 2 \alpha & =\left(a^{\prime}+\alpha^{\prime}\right)-\left(b^{\prime}+\beta^{\prime}\right), \notag\\
2 b^{\prime} & =(a-\alpha)+(b-\beta), & 2 b & =\left(a^{\prime}-\alpha^{\prime}\right)+\left(b^{\prime}-\beta^{\prime}\right), \notag\\
2 \beta^{\prime} & =(a-\alpha)-(b-\beta), & 2 \beta & =\left(a^{\prime}-\alpha^{\prime}\right)-\left(b^{\prime}-\beta^{\prime}\right), \label{eq:8:4:13}
\end{align}
\[
\begin{aligned}
c^{\prime} & =c, & c & =c^{\prime} \\
\gamma^{\prime} & =a-b, & \gamma & =a^{\prime}-b^{\prime}
\end{aligned}
\]
\item
\begin{equation}
\clebsch{a}{\alpha}{b}{\beta}{c}{\gamma}=(-1)^{b+\beta} \sqrt{\frac{2 c+1}{2 c^{\prime}+1}} \clebsch{a^{\prime}}{\alpha^{\prime}}{b^{\prime}}{\beta^{\prime}}{c^{\prime}}{\gamma^{\prime}}, \label{eq:8:4:14}
\end{equation}
where
\begin{align}
2 a^{\prime} & =(b-\beta)+(c+\gamma), & 2 a & =\left(b^{\prime}-\beta^{\prime}\right)+\left(c^{\prime}+\gamma^{\prime}\right), \notag\\
2 \alpha^{\prime} & =2(a+\alpha)-(b-\beta)-(c+\gamma), & 2 \alpha & =2\left(a^{\prime}+\alpha^{\prime}\right)-\left(b^{\prime}-\beta^{\prime}\right)-\left(c^{\prime}+\gamma^{\prime}\right), \notag\\
2 b^{\prime} & =(a-\alpha)+(c+\gamma), & 2 b & =\left(a^{\prime}-\alpha^{\prime}\right)+\left(c^{\prime}+\gamma^{\prime}\right), \label{eq:8:4:15}\\
2 \beta^{\prime} & =2(b+\beta)-(a-\alpha)-(c+\gamma), & 2 \beta & =2\left(b^{\prime}+\beta^{\prime}\right)-\left(a^{\prime}-\alpha^{\prime}\right)-\left(c^{\prime}+\gamma^{\prime}\right), \notag\\
2 c^{\prime} & =(a-\alpha)+(b-\beta), & 2 c & =\left(a^{\prime}-\alpha^{\prime}\right)+\left(b^{\prime}-\beta^{\prime}\right), \notag\\
2 \gamma^{\prime} & =(a-\alpha)+(b-\beta)-2(c-\gamma), & 2 \gamma & =\left(a^{\prime}-\alpha^{\prime}\right)+\left(b^{\prime}-\beta^{\prime}\right)-2\left(c^{\prime}-\gamma^{\prime}\right) . \notag
\end{align}
\end{enumerate}
\subsection{"Mirror" Symmetry}\label{sec:8.4.4}\index{mirror symmetry}\index{symmetry!mirror}\index{negative momenta}
The original relations which define the $3 j m$ symbols and the Clebsch-Gordan coefficients may be extended to the domain of negative integer or half-integer arguments $a, b, c$. In this case the following symmetry relations are valid \cite{ref045} (here we denote $\bar{c} \equiv-c-1, \gamma \equiv-\gamma$, etc.).\\
\subsubsection{$3jm$ symbols}
The relations for the $3 j m$ symbols are
\begin{equation}
\threej{a}{b}{c}{\alpha}{\beta}{\gamma}\label{eq:8:4:16}=(-1)^{b-c-\alpha}\threej{a}{b}{c}{\alpha}{\beta}{\gamma}=i(-1)^{a+\alpha-\beta}\threej{a}{b}{c}{\alpha}{\beta}{\gamma}=i(-1)^{a+b+c}\threej{\bar{a}}{b}{\bar{c}}{\alpha}{\beta}{\gamma} .
\end{equation}
The $3 j \mathrm{~m}$ symbols with negative momenta obey the symmetry relations \eqref{eq:8:4:5} and \eqref{eq:8:4:6} provided $\bar{a}, \bar{b}, \bar{c}$ are replaced by $-a,-b,-c$, respectively, in the phase factors of these relations.\\
\subsubsection{Clebsch-Gordan symbols}
The relations for the Clebsch-Gordan coefficients are
\begin{align}
\clebsch{a}{\alpha}{b}{\beta}{c}{\gamma}=(-1)^{a+b-c} \clebsch{a}{\alpha}{\delta}{\beta}{c}{\gamma} & =(-1)^{b-c-\alpha} \clebsch{a}{\alpha}{b}{\beta}{c}{\gamma}=(-1)^{b+\beta} \clebsch{a}{\alpha}{b}{\beta}{c}{\gamma}=(-1)^{a-\alpha} \clebsch{a}{\alpha}{\delta}{\beta}{c}{\gamma}=\clebsch{a}{\alpha}{b}{\beta}{c}{\gamma}  \notag\\
& =\sqrt{\frac{2 c+1}{2 b+1}} \clebsch{a}{\alpha}{c}{\gamma}{b}{\beta}=\sqrt{\frac{2 c+1}{2 a+1}} \clebsch{c}{\gamma}{b}{\beta}{a}{\alpha}=\sqrt{\frac{2 c+1}{2 a+1}} \clebsch{c}{\gamma}{b}{\beta}{a}{\alpha}. \label{eq:8:4:17}
\end{align}
From \eqref{eq:8:4:17} one gets
\begin{align}
& \clebsch{a}{\bar{a}}{b}{\bar{\beta}}{c}{\bar{\gamma}}=\clebsch{\bar{a}}{\alpha}{\bar{b}}{\beta}{\bar{c}}{\gamma}, \label{eq:8:4:18}\\
& \clebsch{\bar{a}}{\bar{a}}{b}{\beta}{c}{\gamma}=\clebsch{a}{\alpha}{\bar{b}}{\bar{\beta}}{\bar{c}}{\bar{\gamma}}, \label{eq:8:4:19}\\
& \clebsch{a}{\alpha}{\bar{b}}{\bar{\beta}}{c}{\gamma}=\clebsch{\bar{a}}{\bar{a}}{b}{\beta}{\bar{c}}{\bar{\gamma}}. \label{eq:8:4:20}
\end{align}
When tabulating formulas for the Clebsch-Gordan coefficients, it is convenient to use the "mirror" symmetry properties in the form
\begin{equation}
\clebsch{a}{\alpha}{b}{\beta}{a+k}{\gamma}=(-1)^{b+\beta} \clebsch{a}{\alpha}{b}{\beta}{a-k}{\gamma}. \label{eq:8:4:21}
\end{equation}
\subsection{Properties of the Vector-Addition Coefficients under Transformations of the Coordinate System and Time Reversal}\label{sec:8.4.5}\index{vector-addition coefficients!under coordinate transformations}\index{time reversal}\index{coordinate inversion}\isym{Pr-hat@$\widehat{P}_{r}$, coordinate-inversion operator}\isym{Pt-hat@$\widehat{P}_{t}$, time-reversal operator}
The vector-addition coefficients are not invariant under transformations in question, although formally they are independent of coordinates and time.

\subsubsection{Rotation of the coordinate system.}

The projections of angular momenta vary under rotations of the coordinate system. The $3 j m$ symbols and the Clebsch-Gordan coefficients in a new coordinate system may be expressed as superpositions of the corresponding quantities in the initial coordinate system. Thus
\begin{align}
\hat{R}\threej{a}{b}{c}{\alpha}{\beta}{\gamma} & \equiv\threej{a}{b}{c}{\alpha^{\prime}}{\beta^{\prime}}{\gamma^{\prime}}=\sum_{\alpha \beta \gamma}\threej{a}{b}{c}{\alpha}{\beta}{\gamma} D_{\alpha \alpha^{\prime}}^{a}(R) D_{\beta \beta^{\prime}}^{b}(R) D_{\gamma \gamma^{\prime}}^{c}(R), \label{eq:8:4:22}\\
\widehat{R} \clebsch{a}{\alpha}{b}{\beta}{c}{\gamma} & \equiv \clebsch{a}{\alpha^{\prime}}{b}{\beta^{\prime}}{c}{\gamma^{\prime}}=\sum_{\alpha \beta \gamma} \clebsch{a}{\alpha}{b}{\beta}{c}{\gamma} D_{\alpha \alpha^{\prime}}^{a}(R) D_{\beta \beta^{\prime}}^{b}(R) D_{\gamma \gamma^{\prime}}^{c *}(R). \notag
\end{align}
\subsubsection{Inversion of the coordinate system.}

The $3 j m$ symbols and the Clebsch-Gordan coefficients are invariant under inversion of the coordinate system, because angular momentum is a pseudovector.
\begin{align}
\widehat{P}_{r}\threej{a}{b}{c}{\alpha}{\beta}{\gamma} & =\threej{a}{b}{c}{\alpha}{\beta}{\gamma}, \label{eq:8:4:23}\\
\widehat{P}_{r} \clebsch{a}{\alpha}{b}{\beta}{c}{\gamma} & =\clebsch{a}{\alpha}{b}{\beta}{c}{\gamma}. \notag
\end{align}
\subsubsection{(Time reversal.}
The time reversal changes signs of projections of all angular momenta (reversal of the rotation direction). Hence
\begin{align}
\hat{P}_{t}\threej{a}{b}{c}{\alpha}{\beta}{\gamma} & =\threej{a}{b}{c}{-\alpha}{-\beta}{-\gamma}=(-1)^{a+b+c}\threej{a}{b}{c}{\alpha}{\beta}{\gamma}, \label{eq:8:4:24}\\
\widehat{P}_{t} \clebsch{a}{\alpha}{b}{\beta}{c}{\gamma} & =\clebsch{a}{-\alpha}{b}{-\beta}{c}{-\gamma}=(-1)^{a+b-c} \clebsch{a}{\alpha}{b}{\beta}{c}{\gamma}. \notag
\end{align}
\section{EXPLICIT FORMS OF THE CLEBSCH-GORDAN COEFFICIENTS FOR SPECIAL VALUES OF THE ARGUMENTS}\label{sec:8.5}\index{Clebsch-Gordan coefficient!for special values of the arguments}
\subsection{Special Values of Momenta $a, b, c$}\label{sec:8.5.1}\index{Clebsch-Gordan coefficient!for special momenta}
\subsubsection{$c=0$ or $b=0$.}
\begin{align}
& \clebsch{a}{\alpha}{b}{\beta}{0}{0}=(-1)^{a-\alpha} \frac{\delta_{\alpha b} \delta_{\alpha,-\beta}}{\sqrt{2 a+1}}, \label{eq:8:5:1}\\
& \clebsch{a}{\alpha}{0}{0}{\alpha}{\gamma}=\delta_{a c} \delta_{\alpha \gamma}. \label{eq:8:5:2}
\end{align}
\subsubsection{$c=a+b$} 
\begin{equation}
\clebsch{a}{a}{b}{\beta}{a+b}{\alpha+\beta}=\left[\frac{(2 a)!(2 b)!(a+b+\alpha+\beta)!(a+b-\alpha-\beta)!}{(2 a+2 b)!(a+\alpha)!(a-\alpha)!(b+\beta)!(b-\beta)!}\right]^{\frac{1}{2}}. \label{eq:8:5:3}
\end{equation}
The following two formulas are equivalent to \eqref{eq:8:5:3} (i and $n$ are integer)
\begin{gather}
\clebsch{a}{a-n+i}{b}{b-i+1}{a+b}{a+b-n+1}=\left[\frac{\binom{2 a}{n-i}\binom{2 b}{i-1}}{\binom{2 a+2 b}{n-1}}\right]^{\frac{1}{2}}, \label{eq:8:5:4}\\
\clebsch{a}{a+n-i}{b}{-b+i-1}{a+b}{a-b+n-1}=\left[\frac{\binom{2 a}{i-n}\binom{2 b}{i-1}}{\binom{2 a+2 b}{2 a+n-1}}\right]^{\frac{1}{2}}. \label{eq:8:5:5}
\end{gather}
In particular,
\begin{gather}
\clebsch{a}{a}{b}{b}{a+b}{a+b}=1, \label{eq:8:5:6}\\
\clebsch{a}{a}{b}{-b}{a+b}{a-b}=\left[\frac{(2 a)!(2 b)!}{(2 a+2 b)!}\right]^{\frac{1}{2}}. \label{eq:8:5:7}
\end{gather}
\subsubsection{$c=a+b-1$}
\begin{equation}
\clebsch{a}{\alpha}{b}{\beta}{a+b-1}{\alpha+\beta}=2(b \alpha-a \beta)\left[\frac{(2 a+2 b-1)(2 a-1)!(2 b-1)!(a+b+\alpha+\beta-1)!(a+b-\alpha-\beta-1)!}{(a+\alpha)!(a-\alpha)!(b+\beta)!(b-\beta)!(2 a+2 b)!}\right]^{\frac{1}{2}} . \label{eq:8:5:8}
\end{equation}
In particular,
\begin{gather}
\clebsch{a}{\alpha}{b}{\beta}{a+b-1}{\alpha+\beta}=0, \quad \text { if } \quad \frac{\alpha}{\beta}=\frac{a}{b}. \label{eq:8:5:9}\\
\clebsch{a}{a-1}{b}{b}{a+b-1}{a+b-1}=-\sqrt{\frac{b}{a+b}}. \label{eq:8:5:10}\\
\clebsch{a}{a}{b}{b-1}{a+b-1}{a+b-1}=\sqrt{\frac{a}{a+b}}. \label{eq:8:5:11}\\
\clebsch{a}{a}{b}{-b}{a+b-1}{a-b}=\left[\frac{(2 a)!(2 b)!(2 a+2 b-1)}{(2 a+2 b)!}\right]^{\frac{1}{2}}, \label{eq:8:5:12}
\end{gather}
\subsubsection{$c=a-b(a \geq b)$}
\begin{equation}
\clebsch{a}{\alpha}{b}{\beta}{a-b}{\alpha+\beta}=(-1)^{b+\beta}\left[\frac{(a+\alpha)!(a-\alpha)!(2 b)!(2 a-2 b+1)!}{(2 a+1)!(b+\beta)!(b-\beta)!(a-b+\alpha+\beta)!(a-b-\alpha-\beta)!}\right]^{\frac{1}{2}}. \label{eq:8:5:13}
\end{equation}
The following two formulas are equivalent to \eqref{eq:8:5:13} ( $i$ and $n$ are integer)
\begin{gather}
\clebsch{a}{a-n+i}{b}{b-i+1}{a-b}{a+b-n+1}=(-1)^{2 b-i+1}\left[\frac{2 a-2 b+1}{2 a+1} \frac{\binom{2 b}{i-1}\binom{2 a-2 b}{2 a-n+1}}{\binom{2 a}{n-i}}\right]^{\frac{1}{2}}, \label{eq:8:5:14}\\
\clebsch{a}{a-n+i}{b}{-b-i+1}{a-b}{a-b-n+1}=(-1)^{i+1}\left[\frac{2 a-2 b+1}{2 a+1} \frac{\binom{2 b}{-i+1}\binom{2 a-2 b}{n-1}}{\binom{2 a}{n-i}}\right]^{\frac{1}{2}}. \label{eq:8:5:15}
\end{gather}
In particular
\begin{equation}
\clebsch{a}{a}{b}{-b}{a-b}{a-b}=\sqrt{\frac{2 a-2 b+1}{2 a+1}} . \label{eq:8:5:16}
\end{equation}
\subsubsection{$c=a-b+1(a+1 \geq b)$}
\begin{equation}
\clebsch{a}{\alpha}{b}{\beta}{a-b+1}{\alpha+\beta}=(-1)^{b+\beta+1} 2(a \beta+b \alpha+\beta)\left[\frac{(2 a-2 b+3)(2 b-1)!(2 a-2 b+1)!(a+\alpha)!(a-\alpha)!}{(2 a+2)!(b+\beta)!(b-\beta)!(a-b+\alpha+\beta+1)!(a-b-\alpha-\beta+1)!}\right]^{\frac{1}{2}} . \label{eq:8:5:17}
\end{equation}
In particular,
\begin{equation}
\clebsch{a}{a}{b}{-b}{a-b+1}{a-b}=\left[\frac{(2 a-2 b+3) 2 b}{(2 a+2)(2 a+1)}\right]^{\frac{1}{2}}. \label{eq:8:5:18}
\end{equation}
\subsubsection{$c=a+b-2$}
\begin{align}
\clebsch{a}{\alpha}{b}{\beta}{a+b-2}{\alpha+\beta} & =\left[\frac{2 a(2 a-1) 2 b(2 b-1)}{2(2 a+2 b-2)(2 a+2 b-1)}\right]^{\frac{1}{2}}\left[\binom{2 a}{a-\alpha}\binom{2 b}{b-\beta}\binom{2 a+2 b-4}{a+b-\alpha-\beta-2}\right]^{-\frac{1}{2}}  \notag\\
& \times\left\{\binom{2 a-2}{a-\alpha}\binom{2 b-2}{b+\beta}-2\binom{2 a-2}{a-\alpha-1}\binom{2 b-2}{b+\beta-1}+\binom{2 a-2}{a-\alpha-2}\binom{2 b-2}{b+\beta-2}\right\} . \label{eq:8:5:19}
\end{align}
In particular,
\begin{equation}
\clebsch{a}{a}{b}{-b}{a+b-2}{a-b}=\left[\frac{(2 a)!(2 b)!(2 a+2 b-3)}{2(2 a+2 b-1)!}\right]^{\frac{1}{2}}. \label{eq:8:5:20}
\end{equation}
\subsubsection{$c=a-b+2(a+2 \geq b)$}
\begin{gather}
\clebsch{a}{\alpha}{b}{\beta}{a-b+2}{\alpha+\beta}=\left[\frac{(2 b-1) 2 b(2 a-2 b+5)(2 a-2 b+4)(2 a-2 b+3)}{2(2 a+1)(2 a+2)(2 a+3)}\right]^{\frac{1}{2}}  \notag\\
\times\left[\binom{2 a}{a-\alpha}\binom{2 b}{b-\beta}\binom{2 a-2 b+4}{a-b-\alpha-\beta+2}\right]^{-\frac{1}{2}}(-1)^{b+\beta}\left\{\binom{2 b-2}{b+\beta}\binom{2 a-2 b+2}{a-b-\alpha-\beta}\right. \notag\\
\left.-2\binom{2 b-2}{b+\beta-1}\binom{2 a-2 b+2}{a-b-\alpha-\beta+1}+\binom{2 b-2}{b+\beta-2}\binom{2 a-2 b+2}{a-b-\alpha-\beta+2}\right\} . \label{eq:8:5:21}
\end{gather}
In particular,
\begin{equation}
\clebsch{a}{a}{b}{-b}{a-b+2}{a-b}=\left[\frac{(2 a-2 b+5) 2 b(2 b-1)}{(2 a+1)(2 a+2)(2 a+3)}\right]^{\frac{1}{2}}. \label{eq:8:5:22}
\end{equation}
\subsubsection{$a=b, \alpha=\beta$}
\begin{equation}
\clebsch{a}{\alpha}{a}{\alpha}{c}{\gamma}= \begin{cases}0, & \text { if } 2 a+c=2 g+1, \label{eq:8:5:23}\\ \delta_{\gamma, 2 \alpha} \frac{(-1)^{g-c} \sqrt{2 c+1} g!}{\left(\frac{c+\gamma}{2}\right)!\left(\frac{\varepsilon-\gamma}{2}\right)!(g-c)!}\left[\frac{(c+\gamma)!(c-\gamma)!(2 g-2 c)!}{(2 g+1)!}\right]^{\frac{1}{2}}, & \text { if } 2 a+c=2 g,\end{cases}
\end{equation}
$g$ is integer and positive.\\
\subsubsection{$b=a+\frac{1}{2}$, $\beta=\alpha \pm \frac{1}{2}$, $2 a+c=2 g$}
\begin{align}
& \clebsch{a}{\alpha}{a+\frac{1}{2}}{\alpha+\frac{1}{2}}{c+\frac{1}{2}}{2 \alpha+\frac{1}{2}}=\left[\frac{c+2 \alpha+1}{a+\alpha+1}\right]^{\frac{1}{2}} C, \label{eq:8:5:24}\\
& \clebsch{a}{a}{a+\frac{1}{2}}{\alpha-\frac{1}{2}}{c+\frac{1}{2}}{2 \alpha-\frac{1}{2}}=\left[\frac{c-2 \alpha+1}{a-\alpha+1}\right]^{\frac{1}{2}} C, \notag
\end{align}
where
\begin{equation}
C=\frac{(-1)^{c / 2-a}\left(a+\frac{c}{2}\right)!}{\left(a-\frac{c}{2}\right)!\left(\frac{c}{2}+\alpha\right)!\left(\frac{c}{2}-\alpha\right)!}\left[\frac{(2 a+c+2)(2 a-c)!(c+2 \alpha)!(c-2 \alpha)!}{2(2 a+c+1)!}\right]^{\frac{1}{2}}, \label{eq:8:5:25}
\end{equation}
\subsubsection{$b=a+\frac{1}{2}$, $\beta=\alpha \pm \frac{1}{2}$, $2 a+c=2 g+1$}
\begin{align}
& \clebsch{a}{\alpha}{a+\frac{1}{2}}{\alpha+\frac{1}{2}}{c+\frac{1}{2}}{2 \alpha+\frac{1}{2}}=\left[\frac{c-2 \alpha+1}{a+\alpha+1}\right]^{\frac{1}{2}} D, \notag\\
& \clebsch{a}{\alpha}{a+\frac{1}{2}}{\alpha-\frac{1}{2}}{c+\frac{1}{2}}{2 \alpha-\frac{1}{2}}=-\left[\frac{c+2 \alpha+1}{a-\alpha+1}\right]^{\frac{1}{2}} D, \label{eq:8:5:26}
\end{align}
where
\begin{equation}
D=\frac{(-1)^{\frac{a-1}{2}-a}\left(a+\frac{c+1}{2}\right)!}{\left(a-\frac{c+1}{2}\right)!\left(\frac{c+1}{2}+\alpha\right)!\left(\frac{c+1}{2}-\alpha\right)!}\left[\frac{(2 a-c)!(c+2 \alpha+1)!(c-2 \alpha+1)!}{2(2 a+c+2)!}\right]^{\frac{1}{2}}. \label{eq:8:5:27}
\end{equation}
\subsubsection{$b=a+1$; $\beta=\alpha$, $\alpha \pm 1$; $2 a+c=2 g$}
\begin{align}
\clebsch{a}{\alpha}{a+1}{a+1}{c}{2 a+1} & =-\left[\frac{(c+2 \alpha+1)(c-2 \alpha)}{(a+\alpha+2)(a+\alpha+1)}\right]^{\frac{1}{2}} E, \notag\\
\clebsch{a}{\alpha}{a+1}{\alpha}{c}{2 \alpha} & =\frac{2 \alpha}{[(a+\alpha+1)(a-\alpha+1)]^{\frac{1}{2}}} E, \label{eq:8:5:28}\\
\clebsch{a}{\alpha}{a+1}{\alpha-1}{c}{2 \alpha-1} & =\left[\frac{(c-2 \alpha+1)(c+2 \alpha)}{(a-\alpha+2)(a-\alpha+1)}\right]^{\frac{1}{2}} E, \notag
\end{align}
where
\begin{equation}
E=\frac{(-1)^{\frac{c}{2}-a}\left(a+\frac{c}{2}\right)!}{2\left(a-\frac{c}{2}\right)!\left(\frac{c}{2}+\alpha\right)!\left(\frac{c}{2}-\alpha\right)!}\left[\frac{(2 c+1)(2 a+c+2)(2 a-c+1)!(c+2 \alpha)!(c-2 \alpha)!}{c(c+1)(2 a+c+1)!}\right]^{\frac{1}{2}}. \label{eq:8:5:29}
\end{equation}
\subsubsection{$b=a+1$; $\beta=\alpha$, $\alpha \pm 1$; $2 a+c=2 g+1$}
\begin{align}
\clebsch{a}{\alpha}{a+1}{\alpha+1}{c}{2 \alpha+1} & =\frac{c(c+1)+(2 \alpha+1)(2 a+2)}{2\left(\frac{c+1}{2}+\alpha\right)!\left(\frac{c+1}{2}-\alpha\right)!}\left[\frac{(c+2 \alpha+1)!(c-2 \alpha-1)!}{(a+\alpha+2)(a+\alpha+1)}\right]^{\frac{1}{2}} F, \notag\\
\clebsch{a}{\alpha}{a+1}{\alpha}{c}{2 \alpha} & =\frac{2 a+2}{\left(\frac{c-1}{2}+\alpha\right)!\left(\frac{c-1}{2}-\alpha\right)!}\left[\frac{(c+2 \alpha)!(c-2 \alpha)!}{(a+\alpha+1)(a-\alpha+1)}\right]^{\frac{1}{2}} F, \label{eq:8:5:30}\\
\clebsch{a}{\alpha}{a+1}{\alpha-1}{c}{2 \alpha-1} & =\frac{c(c+1)-(2 \alpha-1)(2 a+2)}{2\left(\frac{c+1}{2}-\alpha\right)!\left(\frac{c-1}{2}+\alpha\right)!}\left[\frac{(c+2 \alpha-1)!(c-2 \alpha+1)!}{(a-\alpha+2)(a-\alpha+1)}\right]^{\frac{1}{2}} F, \notag
\end{align}
where
\begin{equation}
F=\frac{(-1)^{\frac{c-1}{2}-a}\left(a+\frac{c+1}{2}\right)!}{\left(a-\frac{c-1}{2}\right)!}\left[\frac{(2 c+1)(2 a-c+1)!}{c(c+1)(2 a+c+2)!}\right]^{\frac{1}{2}}. \label{eq:8:5:31}
\end{equation}
Equations \eqref{eq:8:5:24}-\eqref{eq:8:5:31} were obtained by Stone \cite{ref107}.

\subsection{Special Values of Momentum Projections}\label{sec:8.5.2}\index{Clebsch-Gordan coefficient!for special projections}
\subsubsection{$\alpha=\beta=\gamma=0$}
\begin{equation}
\clebsch{a}{0}{b}{0}{c}{0}= \begin{cases}0, & \text { if } a+b+c=2 g+1  \\ \frac{(-1)^{g-c} \sqrt{2 c+1} g!}{(g-a)!(g-b)!(g-c)!}\left[\frac{(2 g-2 a)!(2 g-2 b)!(2 g-2 c)!}{(2 g+1)!}\right]^{\frac{1}{2}}, & \text { if } a+b+c=2 g\end{cases}, \label{eq:8:5:32}
\end{equation}
$g$ is integer and positive. In particular,
\begin{gather}
\clebsch{a}{0}{b}{0}{a+b}{0}=\frac{(a+b)!}{a!b!}\left[\frac{(2 a)!(2 b)!}{(2 a+2 b)!}\right]^{\frac{1}{2}}, \label{eq:8:5:33}\\
\clebsch{a}{0}{b}{0}{a-b}{0}=(-1)^{b} \frac{a!}{b!(a-b)!}\left[\frac{(2 b)!(2 a-2 b+1)!}{(2 a+1)!}\right]^{\frac{1}{2}}. \label{eq:8:5:34}
\end{gather}
\subsubsection{ $\gamma=c$ or $\alpha=a$ :}
\begin{gather}
\clebsch{a}{\alpha}{b}{\beta}{c}{c}=\delta_{\alpha+\beta, c}(-1)^{a-\alpha}\left[\frac{(2 c+1)!(a+b-c)!(a+\alpha)!(b+\beta)!}{(a+b+c+1)!(a-b+c)!(-a+b+c)!(a-\alpha)!(b-\beta)!}\right]^{\frac{1}{2}}, \label{eq:8:5:35}\\
\clebsch{a}{a}{b}{\beta}{c}{\gamma}=\delta_{\gamma-\beta, \alpha}\left[\frac{(2 c+1)(2 a)!(-a+b+c)!(b-\beta)!(c+\gamma)!}{(a+b+c+1)!(a-b+c)!(a+b-c)!(b+\beta)!(c-\gamma)!}\right]^{\frac{1}{2}}. \label{eq:8:5:36}
\end{gather}
In particular,
\begin{gather}
\clebsch{a}{a}{b}{b}{c}{c}=\delta_{a+b, c}, \label{eq:8:5:37}\\
\clebsch{a}{a}{b}{-b}{c}{c}=\delta_{a-b, c}\left[\frac{2 c+1}{2 a+1}\right]^{\frac{1}{2}}, \label{eq:8:5:38}\\
\clebsch{a}{a}{b}{c-a}{c}{c}=\left[\frac{(2 a)!(2 c+1)!}{(a+b+c+1)!(a-b+c)!}\right]^{\frac{1}{2}}, \label{eq:8:5:39}\\
\clebsch{a}{a-1}{b}{c-a+1}{c}{c}=-\left[\frac{(2 a-1)!(2 c+1)!(a+b-c)(-a+b+c+1)}{(a+b+c+1)!(a-b+c)!}\right]^{\frac{1}{2}}, \label{eq:8:5:40}\\
\clebsch{a}{a}{b}{c-a-1}{c}{c-1}=\left[\frac{(2 a)!(2 c+1)!(a+b-c+1)(-a+b+c)}{(a+b+c+1)!(a-b+c)!2 c}\right]^{\frac{1}{2}}, \label{eq:8:5:41}\\
\clebsch{c}{c}{b}{0}{c}{c}=(2 c)!\left[\frac{2 c+1}{(2 c-b)!(2 c+b+1)!}\right]^{\frac{1}{2}}, \label{eq:8:5:42}\\
\clebsch{c}{c-b}{b}{b}{c}{c}=(-1)^{b}\left[\frac{(2 c+1)!(2 b)!}{(2 c+b+1)!b!}\right]^{\frac{1}{2}}, \label{eq:8:5:43}
\end{gather}
\subsubsection{ $\gamma=c-1$ or $\alpha=a-1$ :}
\begin{align}
\clebsch{a}{\alpha}{b}{\beta}{c}{c-1}= & \delta_{\alpha+\beta, c-1}(-1)^{a-\alpha}\{(b-\beta)(b+\beta+1)-(a-\alpha)(a+\alpha+1)\}  \notag\\
& \times\left[\frac{(2 c+1)(2 c-1)!(a+b-c)!(a+\alpha)!(b+\beta)!}{(a+b+c+1)!(a-b+c)!(-a+b+c)!(a-\alpha)!(b-\beta)!}\right]^{\frac{1}{2}}, \label{eq:8:5:44}\\
\clebsch{a}{a-1}{b}{\beta}{c}{\gamma}= & \delta_{\gamma-\beta, a-1}\{(c-\gamma)(c+\gamma+1)-(b+\beta)(b-\beta+1)\}  \notag\\
& \times\left[\frac{(2 c+1)(2 a-1)!(-a+b+c)!(b-\beta)!(c+\gamma)!}{(a+b+c+1)!(a-b+c)!(a+b-c)!(b+\beta)!(c-\gamma)!}\right]^{\frac{1}{2}}, \label{eq:8:5:45}
\end{align}
In particular,
\begin{gather}
\clebsch{a}{\alpha}{a}{\alpha}{c}{c-1}=0, \label{eq:8:5:46}\\
\clebsch{a}{a-1}{c}{-\gamma}{c}{\gamma}=0, \label{eq:8:5:47}\\
\clebsch{a}{a-1}{b}{\beta}{c}{c-1}=\delta_{\beta, c-a}\{a(a+1)-b(b+1)+c(c+1)-2 a c\}\left[\frac{(2 c+1)(2 c-1)!(2 a-1)!}{(a+b+c+1)!(a-b+c)!}\right]^{\frac{1}{2}} . \label{eq:8:5:48}
\end{gather}
\section{RECURSION RELATIONS FOR THE CLEBSCH-GORDAN COEFFICIENTS}\label{sec:8.6}\index{Clebsch-Gordan coefficient!recursion relations}\index{recursion relations!for the Clebsch-Gordan coefficients}
\subsection{General Recursion Relations}\label{sec:8.6.1}\index{recursion relations!general}
\begin{align}
\clebsch{a}{\alpha}{b}{\beta}{c}{\gamma}= & {\left[\frac{(b+\beta-2 k)!(c+\gamma)!(a+b-c)!(-a+b+c)!(a+b+c+1)!(2 c+1)}{(b+\beta)!(c-\gamma)!(a-b+c)!}\right]^{\frac{1}{2}} }  \notag\\
& \times \sum_{c^{\prime}=c-k}^{c+k}(-1)^{c^{\prime}+k-c} \clebsch{a}{\alpha}{b-k}{\beta-k}{c^{\prime}}{\gamma-k} \frac{\left(c^{\prime}-k+c\right)!(2 k)!}{\left(c+k-c^{\prime}\right)!\left(c+c^{\prime}+k+1\right)!\left(c^{\prime}+k-c\right)!}  \notag\\
& \times\left[\frac{\left(c^{\prime}+k-\gamma\right)!\left(a-b+c^{\prime}+k\right)!\left(2 c^{\prime}+1\right)}{\left(c^{\prime}-k+\gamma\right)!\left(-a+b+c^{\prime}-k\right)!\left(a+b-c^{\prime}-k\right)!\left(a+b+c^{\prime}-k+1\right)!}\right]^{\frac{1}{2}}. \label{eq:8:6:1}
\end{align}
(Yutsis and Bandzaitis \cite{ref045}) where $k$ is integer or half-integer, $0 \leq k \leq(b+\beta) / 2$. In accordance with the choice of $k$ \eqref{eq:8:6:1} yields different recursion relations.

Note also the relation obtained by Stone \cite{ref107}
\begin{align}
\clebsch{a}{\alpha}{b}{\beta}{c}{\gamma}= & {\left[\frac{(a+b+c+1)!(a-b+c)!(a+b-c)!(b+\beta)!(b-\beta)!(a-b+\alpha-\beta)!(a-b-\alpha+\beta)!}{(-a+b+c)!(a+\alpha)!(a-\alpha)!}\right]^{\frac{1}{2}} }  \notag\\
& \times \sum_{b^{\prime}}(-1)^{b-b^{\prime} / 2} 
\clebsch{a-b}{\alpha-\beta}{b^{\prime}}{2 \beta}{c}{\gamma} \frac{\left(2 b^{\prime}+1\right)\left(b+\frac{b^{\prime}}{2}\right)!}{\left(2 b+b^{\prime}+1\right)!\left(b-\frac{b^{\prime}}{2}\right)!\left(\frac{b^{\prime}}{2}+\beta\right)!\left(\frac{b^{\prime}}{2}-\beta\right)!}  \notag\\
& \times\left[\frac{\left(-a+b+c+b^{\prime}\right)!\left(b^{\prime}+2 \beta\right)!\left(b^{\prime}-2 \beta\right)!}{\left(a-b+c+b^{\prime}+1\right)!\left(a-b+c-b^{\prime}\right)!\left(a-b-c+b^{\prime}\right)!}\right]^{\frac{1}{2}}. \label{eq:8:6:2}
\end{align}
Here $a-b \geq|\alpha-\beta| \geq 0, b^{\prime}$ is integer, $b^{\prime}+2 b$ is even, and $a-b+c \geq b^{\prime} \geq|-a+b+c|, 2 b \geq b^{\prime} \geq|2 \beta|$.\\[0pt]
Equation \eqref{eq:8:6:1} may be rewritten by the use of quasi-powers in the form \cite{ref045}
\begin{align}
\clebsch{a}{\alpha}{b}{\beta}{c}{\gamma}= & { \left.\left[\frac{2 c+1}{(b+\beta)^{(2 k)}}\right]^{\frac{1}{2}} \sum_{k^{\prime}=-k}^{k}(-1)^{k+k^{\prime}} \clebsch{a}{\alpha}{b-k}{\beta-k}{c+k^{\prime}}{\gamma-k} \right\rvert\,(c+\gamma)^{\left(k-k^{\prime}\right)}\left(c-\gamma+k+k^{\prime}\right)^{\left(k+k^{\prime}\right)}(a+b-c)^{\left(k+k^{\prime}\right)} }  \notag\\
& \left.\times(-a+b+c)^{\left(k-k^{\prime}\right)}(a+b+c+1)^{\left(k-k^{\prime}\right)}\left(a-b+c+k+k^{\prime}\right)^{\left(k+k^{\prime}\right)}\left(2 c+2 k^{\prime}+1\right)\right]^{\frac{1}{2}}  \notag\\
& \times \frac{(2 k)^{\left(k+k^{\prime}\right)}}{\left(2 c+k+k^{\prime}+1\right)^{(2 k+1)}\left(k+k^{\prime}\right)^{\left(k+k^{\prime}\right)}}, \label{eq:8:6:3}
\end{align}
where
\[
a^{(n)}=\frac{a!}{(a-n)!}, \quad \frac{1}{2} \leq k \leq \frac{b-\kappa}{2}, \quad \kappa= \begin{cases}0 & \text { if } b \text { is integer }, \\ \frac{1}{2} & \text { if } b \text { is half-integer } .\end{cases}
\]
\subsection{Arguments $\alpha, \beta, \gamma$ Change by 1}\label{sec:8.6.2}\index{recursion relations!in the projections}
\begin{equation}
[(c \pm \gamma)(c \mp \gamma+1)]^{\frac{1}{2}} \clebsch{a}{\alpha}{b}{\beta}{c}{\gamma \mp 1}=[(a \mp \alpha)(a \pm \alpha+1)]^{\frac{1}{2}} \clebsch{a}{\alpha \pm 1}{b}{\beta}{c}{\gamma}+[(b \mp \beta)(b \pm \beta+1)]^{\frac{1}{2}} \clebsch{a}{\alpha}{b}{\beta \pm 1}{c}{\gamma}. \label{eq:8:6:4}
\end{equation}
In particular, for $|\gamma|=c$ one has
\begin{equation}
\clebsch{a}{\alpha \mp 1}{ b}{\beta}{c}{ \pm c}=-
\clebsch{a}{\alpha}{ b}{\beta\mp 1}{c}{\pm c}\left[\frac{(b \pm \beta)(b \mp \beta+1)}{(a \pm \alpha)(a \mp \alpha+1)}\right]^{\frac{1}{2}}, \label{eq:8:6:5}
\end{equation}
when $|\alpha|=|\beta|=\frac{1}{2}$
\begin{equation}
\clebsch{a}{\frac{1}{2}}{b}{\frac{1}{2}}{c}{1}=\clebsch{a}{\frac{1}{2}}{b}{-\frac{1}{2}}{c}{0} \frac{(2 b+1)+(-1)^{a+b-c}(2 a+1)}{2[c(c+1)]^{\frac{1}{2}}}, \label{eq:8:6:6}
\end{equation}
and for $|\alpha|=|\beta|=1, a+b+c$ even, one gets
\begin{gather}
\clebsch{a}{1}{b}{-1}{c}{0}=\clebsch{a}{0}{b}{0}{c}{0} \frac{c(c+1)-a(a+1)-b(b+1)}{2[a(a+1) b(b+1)]^{\frac{1}{2}}}, \label{eq:8:6:7}\\
\clebsch{a}{1}{b}{1}{c}{2}=\clebsch{a}{0}{b}{0}{c}{0} \frac{a(a+1)[c(c+1)-a(a+1)+b(b+1)]+b(b+1)[c(c+1)+a(a+1)-b(b+1)]}{2[a(a+1) b(b+1)(c-1) c(c+1)(c+2)]^{\frac{1}{2}}}. \label{eq:8:6:8}
\end{gather}
\subsection{Arguments Change by $1/2$}\label{sec:8.6.3}\index{recursion relations!in half-integer steps}
\begin{gather}
(2 a+1)[b \pm \beta]^{\frac{1}{2}} \clebsch{a}{\alpha}{b}{\beta}{c}{\gamma}=\mp[(a \mp \alpha)(a+b-c)(a+b+c+1)]^{\frac{1}{2}} \clebsch{a-\frac{1}{2}}{\alpha \pm \frac{1}{2}}{b-\frac{1}{2}}{\beta \mp \frac{1}{2}}{c}{\gamma}  \notag\\
+[(a \pm \alpha+1)(-a+b+c)(a-b+c+1)]^{\frac{1}{2}} \clebsch{a+\frac{1}{2}}{\alpha \pm \frac{1}{2}}{b-\frac{1}{2}}{\beta \mp \frac{1}{2}}{c}{\gamma}  \label{eq:8:6:9}
\end{gather}
\begin{gather}
(2 a+1)[b \mp \beta+1]^{\frac{1}{2}} \clebsch{a}{\alpha}{b}{\beta}{c}{\gamma}=[(a \mp \alpha)(a-b+c)(-a+b+c+1)]^{\frac{1}{2}} \clebsch{a-\frac{1}{2}}{\alpha \pm \frac{1}{2}}{b+\frac{1}{2}}{\beta \mp \frac{1}{2}}{c}{\gamma}  \notag\\
\pm[(a \pm \alpha+1)(a+b-c+1)(a+b+c+2)]^{\frac{1}{2}} \clebsch{a+\frac{1}{2}}{\alpha \pm \frac{1}{2}}{b+\frac{1}{2}}{\beta \mp \frac{1}{2}}{c}{\gamma}  \label{eq:8:6:10}
\end{gather}
\begin{multline}
{[(-a+b+c)(a-b+c+1)]^{\frac{1}{2}} \clebsch{a}{\alpha}{b}{\beta}{c}{\gamma}=[(a-\alpha+1)(b-\beta)]^{\frac{1}{2}} \clebsch{a+\frac{1}{2}}{\alpha-\frac{1}{2}}{b-\frac{1}{2}}{\beta+\frac{1}{2}}{c}{\gamma}}  \\
+[(a+\alpha+1)(b+\beta)]^{\frac{1}{2}} \clebsch{a+\frac{1}{2}}{\alpha+\frac{1}{2}}{b-\frac{1}{2}}{\beta-\frac{1}{2}}{c}{\gamma}  \label{eq:8:6:11}
\end{multline}
\begin{multline}
{[(a-b+c)(-a+b+c+1)]^{\frac{1}{2}} \clebsch{a}{\alpha}{b}{\beta}{c}{\gamma}=[(a-\alpha)(b-\beta+1)]^{\frac{1}{2}} \clebsch{a-\frac{1}{2}}{\alpha+\frac{1}{2}}{b+\frac{1}{2}}{\beta-\frac{1}{2}}{c}{\gamma}}  \\
+[(a+\alpha)(b+\beta+1)]^{\frac{1}{2}} \clebsch{a-\frac{1}{2}}{\alpha-\frac{1}{2}}{b+\frac{1}{2}}{\beta+\frac{1}{2}}{c}{\gamma}  \label{eq:8:6:12}
\end{multline}
\begin{multline}
\left[\frac{2 c(-a+b+c)(a+b+c+1)}{2 c+1}\right]^{\frac{1}{2}}
\clebsch{a}{\alpha}{b}{\beta}{c}{\gamma}=
[(b-\beta)(c-\gamma)]^{\frac{1}{2}} 
\clebsch{a}{\alpha}{b-\frac{1}{2}}{\beta+\frac{1}{2}}{c-\frac{1}{2}}{\gamma+\frac{1}{2}}
+[(b+\beta)(c+\gamma)]^\frac{1}{2} \clebsch{a}{\alpha}{b-\frac{1}{2}}{\beta-\frac{1}{2}}{c-\frac{1}{2}}{\gamma-\frac{1}{2}} , \label{eq:8:6:13}
\end{multline}
\begin{multline}
{\left[\frac{2 c(c \mp \gamma)(a+b+c+1)}{2 c+1}\right]}^{\frac{1}{2}}
\clebsch{a}{\alpha}{b}{\beta}{c}{\gamma}
=[(a \mp \alpha)(a-b+c)]^{\frac{1}{2}} 
\clebsch{a-\frac{1}{2}}{\alpha\pm\frac{1}{2}}{b}{\beta}{c-\frac{1}{2}}{\gamma\pm\frac{1}{2}}+
[(b\mp\beta)(-a+b+c)]^{\frac{1}{2}} 
\clebsch{a}{\alpha}{b-\frac{1}{2}}{\beta\pm\frac{1}{2}} {c-\frac{1}{2}}{\gamma\pm\frac{1}{2}}, \label{eq:8:6:14}
\end{multline}
\begin{multline}
  {\left[\frac{2 c(c \mp \gamma)(a+b-c+1)}{2 c+1}\right]}^{\frac{1}{2}}
\clebsch{a}{\alpha}{b}{\beta}{c}{\gamma}
=\pm[(a \pm \alpha+1)(-a+b+c)]^{\frac{1}{2}} 
\clebsch{a+\frac{1}{2}}{\alpha\pm\frac{1}{2}}{b}{\beta}{c-\frac{1}{2}}{\gamma\pm\frac{1}{2}}\\
\mp[(b \pm \beta+1)(a-b+c)]^{\frac{1}{2}} \clebsch{a}{\alpha}{b+\frac{1}{2}}{\beta \pm \frac{1}{2}}{c-\frac{1}{2}}{\gamma \pm \frac{1}{2}} , \label{eq:8:6:15}
\end{multline}
\begin{multline}
{\left[\frac{2(c+1)(c \pm \gamma+1)(a+b-c)}{2 c+1}\right]^{\frac{1}{2}}\clebsch{a}{\alpha}{b}{\beta}{c}{\gamma}
=\mp[(a \mp \alpha)(-a+b+c+1)]^{\frac{1}{2}} \clebsch{a-\frac{1}{2}}{\alpha \pm \frac{1}{2}}{b}{\beta}{c+\frac{1}{2}}{\gamma \pm \frac{1}{2}}} \\
\pm[(b \mp \beta)(a-b+c+1)]^{\frac{1}{2}} \clebsch{a}{\alpha}{b-\frac{1}{2}}{\beta \pm \frac{1}{2}}{c+\frac{1}{2}}{\gamma \pm \frac{1}{2}}, \label{eq:8:6:16}
\end{multline}
\begin{multline}
{\left[(2c+1)(b\mp\beta)\right]^{\frac{1}{2}}\clebsch{a}{\alpha}{b}{\beta}{c}{\gamma}=\left[\frac{(c\mp\gamma)(-a+b+c)(a+b+c+1)}{2c}\right]^{\frac{1}{2}} \clebsch{a}{\alpha}{b-\frac{1}{2}}{\beta \pm \frac{1}{2}}{c-\frac{1}{2}}{\gamma \pm \frac{1}{2}}} \\
\pm\left[\frac{(c \pm \gamma+1)(a-b+c+1)(a+b-c)}{2(c+1)}\right]^{\frac{1}{2}}
\clebsch{a}{\alpha}{b-\frac{1}{2}}{\beta \pm \frac{1}{2}}{c+\frac{1}{2}}{\gamma\pm\frac{1}{2}}. 
\label{eq:8:6:17}
\end{multline}
In particular, Eqs.~\eqref{eq:8:6:9} and \eqref{eq:8:6:10} yield
\begin{equation}
\clebsch{a}{\pm \frac{1}{2}}{b}{\mp \frac{1}{2}}{c}{0}= \pm \clebsch{a-\frac{1}{2}}{0}{b-\frac{1}{2}}{0}{c}{0}\left[\frac{(a+b-c)(a+b+c+1)}{(2 a+1)(2 b+1)}\right]^{\frac{1}{2}}=\mp \clebsch{a+\frac{1}{2}}{0}{b+\frac{1}{2}}{0}{c}{0}\left[\frac{(a+b-c+1)(a+b+c+2)}{(2 a+1)(2 b+1)}\right]^{\frac{1}{2}}, \label{eq:8:6:18}
\end{equation}
if $a+b+c$ is odd, and
\begin{equation}
\clebsch{a}{\pm \frac{1}{2}}{b}{\mp \frac{1}{2}}{c}{0}=\clebsch{a+\frac{1}{2}}{0}{b-\frac{1}{2}}{0}{c}{0}\left[\frac{(-a+b+c)(a-b+c+1)}{(2 a+1)(2 b+1)}\right]^{\frac{1}{2}}=\clebsch{a-\frac{1}{2}}{0}{b+\frac{1}{2}}{0}{c}{0}\left[\frac{(a-b+c)(-a+b+c+1)}{(2 a+1)(2 b+1)}\right]^{\frac{1}{2}}, \label{eq:8:6:19}
\end{equation}
if $a+b+c$ is even.

\subsection{The Case When $\alpha=\beta=\gamma=0$}\label{sec:8.6.4}\index{recursion relations!for zero projections}
Let $p$ be integer and $2 g=a+b+c$. Then
\begin{equation}
\clebsch{a+p}{0}{b-p}{0}{c}{0}=\clebsch{a}{0}{b}{0}{c}{0} \frac{(g-a)!(g-b)!}{(g-a-p)!(g-b+p)!}\left[\frac{(2 g-2 a-2 p)!(2 g-2 b+2 p)!}{(2 g-2 a)!(2 g-2 b)!}\right]^{\frac{1}{2}}. \label{eq:8:6:20}
\end{equation}
In particular, \eqref{eq:8:6:20} yields
\begin{gather}
\clebsch{a+1}{0}{b-1}{0}{c}{0}=\clebsch{a}{0}{b}{0}{c}{0}\left[\frac{(-a+b+c)(a-b+c+1)}{(-a+b+c-1)(a-b+c+2)}\right]^{\frac{1}{2}}, \label{eq:8:6:21}\\
\clebsch{a}{0}{b+2p}{0}{c}{0}=\clebsch{a}{0}{b}{0}{c}{0}(-1)^{p} \frac{(g+p)!(g-a)!(g-b)!(g-c)!}{g!(g-a+p)!(g-b-p)!(g-c+p)!}, \notag\\
\times\left[\frac{(a+b-c+2 p)!(a-b+c-2 p)!(-a+b+c+2 p)!(a+b+c+1)!}{(a+b-c)!(a-b+c)!(-a+b+c)!(a+b+c+2 p+1)!}\right]^{\frac{1}{2}}. \label{eq:8:6:22}
\end{gather}
Moreover, from \eqref{eq:8:6:22} one gets
\begin{equation}
\clebsch{a}{0}{b+2}{0}{c}{0}=-\clebsch{a}{0}{b}{0}{c}{0}\left[\frac{(a+b+c+2)(a+b-c+1)(a-b+c)(-a+b+c+1)}{(a+b+c+3)(a+b-c+2)(a-b+c-1)(-a+b+c+2)}\right]^{\frac{1}{2}}. \label{eq:8:6:23}
\end{equation}
\subsection{Arguments $a, b, c$ Change by 1}\label{sec:8.6.5}\index{recursion relations!in the momenta}
\begin{align}
& 2\left[b^{2}-\beta^{2}\right]^{\frac{1}{2}} \clebsch{a}{\alpha}{b-1}{\beta}{c}{\gamma}=\frac{1}{a(2 a+1)}\left[\left(a^{2}-\alpha^{2}\right)(-a+b+c)(-a+b+c+1)(a-b+c)(a-b+c+1)\right]^{\frac{1}{2}} \notag\\
& \quad \times \clebsch{a-1}{\alpha}{b}{\beta}{c}{\gamma}+\frac{\alpha}{a(a+1)}[(-a+b+c)(a-b+c+1)(a+b-c)(a+b+c+1)]^{\frac{1}{2}} \clebsch{a}{\alpha}{b}{\beta}{c}{\gamma}  \notag\\
& -\frac{1}{(a+1)(2 a+1)}\left\{\left[(a+1)^{2}-\alpha^{2}\right](a+b-c)(a+b-c+1)(a+b+c+1)(a+b+c+2)\right\}^{\frac{1}{2}} \clebsch{a+1}{\alpha}{b}{\beta}{c}{\gamma}, \label{eq:8:6:24}\\
& 2 \beta \clebsch{a}{\alpha}{b}{\beta}{c}{\gamma}=-\frac{1}{a(2 a+1)}\left[\left(a^{2}-\alpha^{2}\right)(-a+b+c+1)(a-b+c)(a+b-c)(a+b+c+1)\right]^{\frac{1}{2}} \clebsch{a-1}{\alpha}{b}{\beta}{c}{\gamma}  \notag\\
& \quad-\frac{\alpha}{a(a+1)}\{a(a+1)+b(b+1)-c(c+1)\} \clebsch{a}{\alpha}{b}{\beta}{c}{\gamma}  \notag\\
& -\frac{1}{(a+1)(2 a+1)}\left\{\left[(a+1)^{2}-\alpha^{2}\right](-a+b+c)(a-b+c+1)(a+b-c+1)(a+b+c+2)\right\}^{\frac{1}{2}} \clebsch{a+1}{\alpha}{b}{\beta}{c}{\gamma}, \label{eq:8:6:25}
\end{align}
\begin{gather}
2\left[(b+1)^{2}-\beta^{2}\right]^{\frac{1}{2}} \clebsch{a}{\alpha}{b+1}{\beta}{c}{\gamma}=-\frac{1}{a(2 a+1)}\left[\left(a^{2}-\alpha^{2}\right)(a+b-c)(a+b-c+1)(a+b+c+1)(a+b+c+2)\right]^{\frac{1}{2}} \notag\\
\times \clebsch{a-1}{\alpha}{b}{\beta}{c}{\gamma}+\frac{\alpha}{a(a+1)}[(-a+b+c+1)(a-b+c)(a+b-c+1)(a+b+c+2)]^{\frac{1}{2}} \clebsch{a}{\alpha}{b}{\beta}{c}{\gamma}  \notag\\
+\frac{1}{(a+1)(2 a+1)}\left\{\left[(a+1)^{2}-\alpha^{2}\right](-a+b+c)(-a+b+c+1)(a-b+c)(a-b+c+1)\right\}^{\frac{1}{2}} \clebsch{a+1}{\alpha}{b}{\beta}{c}{\gamma}, \label{eq:8:6:26}\\
\clebsch{a}{\alpha}{b}{\beta}{c}{\gamma}=\left[\frac{4 c^{2}(2 c+1)(2 c-1)}{(c+\gamma)(c-\gamma)(-a+b+c)(a-b+c)(a+b-c+1)(a+b+c+1)}\right]^{\frac{1}{2}}  \notag\\
\times\left\{\frac{(\alpha-\beta) c(c-1)-\gamma a(a+1)+\gamma b(b+1)}{2 c(c-1)} \clebsch{a}{\alpha}{b}{\beta}{c-1}{\gamma}\right. \notag\\
\left.-\left[\frac{(c-\gamma-1)(c+\gamma-1)(-a+b+c-1)(a-b+c-1)(a+b-c+2)(a+b+c)}{4(c-1)^{2}(2 c-3)(2 c-1)}\right]^{\frac{1}{2}} \clebsch{a}{\alpha}{b}{\beta}{c-2}{\gamma}\right\} . \label{eq:8:6:27}
\end{gather}
\subsection{Arguments $a, b, \alpha, \beta$ Change by 1}\label{sec:8.6.6}\index{recursion relations!in $a, b, \alpha, \beta$}
\begin{gather}
{[(b \pm \beta)(b \pm \beta+1)]^{\frac{1}{2}} \clebsch{a}{\alpha}{b-1}{\beta}{c}{\gamma}}  \notag\\
=\frac{1}{2 a(2 a+1)}[(a \pm \alpha)(a \pm \alpha-1)(-a+b+c)(-a+b+c+1)(a-b+c)(a-b+c+1)]^{\frac{1}{2}} \clebsch{a-1}{\alpha \mp 1}{b}{\beta \pm 1}{c}{\gamma}  \notag\\
\mp \frac{1}{2 a(a+1)}[(a \pm \alpha)(a \mp \alpha+1)(-a+b+c)(a-b+c+1)(a+b-c)(a+b+c+1)]^{\frac{1}{2}} \clebsch{a}{\alpha \mp 1}{b}{\beta \pm 1}{c}{\gamma}  \notag\\
+\frac{1}{2(a+1)(2 a+1)}[(a \mp \alpha+1)(a \mp \alpha+2)(a+b-c)(a+b-c+1)(a+b+c+1)(a+b+c+2)]^{\frac{1}{2}}  \notag\\
\times \clebsch{a+1}{\alpha \mp 1}{b}{\beta \pm 1}{c}{\gamma} . \label{eq:8:6:28}\end{gather}
\begin{gather}
[(b\mp\beta)(b\pm\beta+1)]^{\frac{1}{2}}\clebsch{a}{\alpha}{b}{\beta}{c}{\gamma}\notag\\
= \pm \frac{1}{2 a(2 a+1)}[(a \pm \alpha)(a \pm \alpha-1)(-a+b+c+1)(a-b+c)(a+b-c)(a+b+c+1)]^{\frac{1}{2}} \clebsch{a-1}{\alpha \mp 1}{b}{\beta \pm 1}{c}{\gamma}  \notag\\
\quad-\frac{1}{2 a(a+1)}\{a(a+1)+b(b+1)-c(c+1)\}\left[(a \pm \alpha)(a \mp \alpha+1)\right]^{\frac{1}{2}} \clebsch{a}{\alpha \mp 1}{b}{\beta \pm 1}{c}{\gamma}  \notag\\
\mp \frac{1}{2(a+1)(2 a+1)}[(a \mp \alpha+1)(a \mp \alpha+2)(-a+b+c)(a-b+c+1)(a+b-c+1)(a+b+c+2)]^{\frac{1}{2}}  \notag\\
\quad \times \clebsch{a+1}{\alpha \mp 1}{b}{\beta \pm 1}{c}{\gamma} . \label{eq:8:6:29}
\end{gather}
\begin{gather}
{[(b \mp \beta)(b \mp \beta+1)]^{\frac{1}{2}} \clebsch{a}{\alpha}{b+1}{\beta}{c}{\gamma}}  \notag\\
=\frac{1}{2 a(2 a+1)}[(a \pm \alpha)(a \pm \alpha-1)(a+b-c)(a+b-c+1)(a+b+c+1)(a+b+c+2)]^{\frac{1}{2}} \clebsch{a-1}{\alpha \mp 1}{b}{\beta \pm 1}{c}{\gamma}  \notag\\
\pm \frac{1}{2 a(a+1)}[(a \pm \alpha)(a \mp \alpha+1)(-a+b+c+1)(a-b+c)(a+b-c+1)(a+b+c+2)]^{\frac{1}{2}} \clebsch{a}{\alpha \mp 1}{b}{\beta \pm 1}{c}{\gamma}  \notag\\
+\frac{1}{2(a+1)(2 a+1)}[(a \mp \alpha+1)(a \mp \alpha+2)(-a+b+c)(-a+b+c+1)(a-b+c)(a-b+c+1)]^{\frac{1}{2}}  \notag\\
\times \clebsch{a+1}{\alpha \mp 1}{b}{\beta \pm 1}{c}{\gamma} \label{eq:8:6:30}
\end{gather}
\subsection{Arguments $c, b, \gamma, \beta$ Change by 1}\label{sec:8.6.7}\index{recursion relations!in $c, b, \gamma, \beta$}
\begin{gather}
{\left[\frac{(-a+b+c)(a-b+c)(a+b-c+1)(a+b+c+1)(2 c-1)}{(2 c+1)}\right]^{\frac{1}{2}} \clebsch{a}{\alpha}{b}{\beta}{c}{\gamma}}  \notag\\
=[(b+\beta)(b-\beta+1)(c+\gamma)(c+\gamma-1)]^{\frac{1}{2}} \clebsch{a}{\alpha}{b}{\beta-1}{c-1}{\gamma-1}-2 \beta\left[c^{2}-\gamma^{2}\right]^{\frac{1}{2}} \clebsch{a}{\alpha}{b}{\beta}{c-1}{\gamma}  \notag\\
-[(b-\beta)(b+\beta+1)(c-\gamma)(c-\gamma-1)]^{\frac{1}{2}} \clebsch{a}{\alpha}{b}{\beta+1}{c-1}{\gamma+1} . \label{eq:8:6:31}
\end{gather}
\begin{gather}
{[(b \mp \beta)(b \mp \beta+1)]^{\frac{1}{2}} \clebsch{a}{\alpha}{b-1}{\beta}{c}{\gamma}}  \notag\\
=\frac{1}{2 c}\left[\frac{(c \pm \gamma)(c \pm \gamma-1)(a-b+c)(a-b+c+1)(a+b-c)(a+b-c+1)}{(2 c-1)(2 c+1)}\right]^{\frac{1}{2}} \clebsch{a}{\alpha}{b}{\beta \mp 1}{c-1}{\gamma \mp 1}  \notag\\
\pm \frac{1}{2 c(c+1)}[(c \pm \gamma)(c \mp \gamma+1)(-a+b+c)(a-b+c+1)(a+b-c)(a+b+c+1)]^{\frac{1}{2}} \clebsch{a}{\alpha}{b}{\beta \mp 1}{c}{\gamma \mp 1}  \notag\\
+\frac{1}{2(c+1)}\left[\frac{(c \mp \gamma+1)(c \mp \gamma+2)(-a+b+c)(-a+b+c+1)(a+b+c+1)(a+b+c+2)}{(2 c+1)(2 c+3)}\right]^{\frac{1}{2}} \clebsch{a}{\alpha}{b}{\beta \mp 1}{c+1}{\gamma \mp 1}  \label{eq:8:6:32}
\end{gather}
\begin{gather}
  [(b\pm\beta)(b\mp\beta+1)]^{\frac{1}{2}}\clebsch{a}{\alpha}{b}{\beta}{c}{\gamma}\notag\\
= \pm \frac{1}{2 c}\left[\frac{(c \pm \gamma)(c \pm \gamma-1)(-a+b+c)(a-b+c)(a+b-c+1)(a+b+c+1)}{(2 c-1)(2 c+1)}\right]^{\frac{1}{2}} \clebsch{a}{\alpha}{b}{\beta \mp 1}{c-1}{\gamma \mp 1}  \notag\\
+\frac{1}{2 c(c+1)}\{-a(a+1)+b(b+1)+c(c+1)\}[(c \pm \gamma)(c \mp \gamma+1)]^{\frac{1}{2}} \clebsch{a}{\alpha}{b}{\beta \mp 1}{c}{\gamma \mp 1}  \notag\\
\mp \frac{1}{2(c+1)}\left[\frac{(c \mp \gamma+1)(c \mp \gamma+2)(-a+b+c+1)(a-b+c+1)(a+b-c)(a+b+c+2)}{(2 c+1)(2 c+3)}\right]^{\frac{1}{2}} \clebsch{a}{\alpha}{b}{\beta \mp 1}{c+1}{\gamma \mp 1} \label{eq:8:6:33}
\end{gather}
\begin{gather}
  [(b\pm\beta)(b\pm\beta+1)]^{\frac{1}{2}}\clebsch{a}{\alpha}{b+1}{\beta}{c}{\gamma}\notag\\
=\frac{1}{2 c}\left[\frac{(c \pm \gamma)(c \pm \gamma-1)(-a+b+c)(-a+b+c+1)(a+b+c+1)(a+b+c+2)}{(2 c-1)(2 c+1)}\right]^{\frac{1}{2}} \clebsch{a}{\alpha}{b}{\beta \mp 1}{c-1}{\gamma \mp 1} \notag\\
\mp \frac{1}{2 c(c+1)}[(c \pm \gamma)(c \mp \gamma+1)(-a+b+c+1)(a-b+c)(a+b-c+1)(a+b+c+2)]^{\frac{1}{2}} \clebsch{a}{\alpha}{b}{\beta \mp 1}{c}{\gamma \mp 1}  \notag\\
+\frac{1}{2(c+1)}\left[\frac{(c \mp \gamma+1)(c \mp \gamma+2)(a-b+c)(a-b+c+1)(a+b-c)(a+b-c+1)}{(2 c+1)(2 c+3)}\right]^{\frac{1}{2}} \clebsch{a}{\alpha}{b}{\beta \mp 1}{c+1}{\gamma \mp 1} .  \label{eq:8:6:34}
\end{gather}
\subsection{Recursion Relations for the Regge Symbols}\label{sec:8.6.8}\index{Regge R-symbol@Regge $R$-symbol!recursion relations}
Recursion relations for the Regge $R$-symbols are very convenient because each relation is equivalent to a variety of formally different expressions for the Clebsch-Gordan coefficients. The formulas given below represent some recursion relations $[86,103]$. In these formulas we use the notation $J=\sum_{k=1}^{3} R_{i k}=\sum_{i=1}^{3} R_{i k}=a+b+c$.
\begin{align}
& {\left[R_{11}(J+1)\right]^{\frac{1}{2}}\left\|\begin{array}{lll}
R_{11} & R_{12} & R_{13} \\
R_{21} & R_{22} & R_{23} \\
R_{31} & R_{32} & R_{33}
\end{array}\right\|-\left[R_{22} R_{33}\right]^{\frac{1}{2}}\left\|\begin{array}{lll}
R_{11}-1 & R_{12} & R_{13} \\
R_{21} & R_{22}-1 & R_{23} \\
R_{31} & R_{32} & R_{33}-1
\end{array}\right\|}  \notag\\
& +\left[R_{23} R_{32}\right]^{\frac{1}{2}}\left\|\begin{array}{lll}
R_{11}-1 & R_{12} & R_{13} \\
R_{21} & R_{22} & R_{23}-1 \\
R_{31} & R_{32}-1 & R_{33}
\end{array}\right\|=0 . \label{eq:8:6:35}\\
& {\left[\left(R_{11}+1\right) R_{21}\right]^{\frac{1}{2}}\left\|\begin{array}{lll}
R_{11}+1 & R_{12} & R_{13} \\
R_{21}-1 & R_{22} & R_{23} \\
R_{31} & R_{32} & R_{33}
\end{array}\right\|+\left[\left(R_{12}+1\right) R_{22}\right]^{\frac{1}{2}}\left\|\begin{array}{lll}
R_{11} & R_{12}+1 & R_{13} \\
R_{21} & R_{22}-1 & R_{23} \\
R_{31} & R_{32} & R_{33}
\end{array}\right\|}  \notag\\
& +\left[\left(R_{13}+1\right) R_{23}\right]^{\frac{1}{2}}\left\|\begin{array}{lll}
R_{11} & R_{12} & R_{13}+1 \\
R_{21} & R_{22} & R_{23}-1 \\
R_{31} & R_{32} & R_{33}
\end{array}\right\|=0 .
\footnote{Column recursion — it shifts R₁ₖ+1 and R₂ₖ−1, so each array has equal column sums but row sums J+1, J−1, J. A Regge symbol requires all six line-sums equal, so every term vanishes as a standard 3jm and the relation is only vacuously satisfied.  It's not an OCR issue; the transcription matches the scan.}
\label{eq:8:6:36}\\
& {\left[R_{11} R_{22}(J+1)\right]^{\frac{1}{2}}\left\|\begin{array}{lll}
R_{11} & R_{12} & R_{13} \\
R_{21} & R_{22} & R_{23} \\
R_{31} & R_{32} & R_{33}
\end{array}\right\|-\left(R_{22}+R_{23}\right)\left[R_{33}\right]^{\frac{1}{2}}\left\|\begin{array}{lll}
R_{11}-1 & R_{12} & R_{13} \\
R_{21} & R_{22}-1 & R_{23} \\
R_{31} & R_{32} & R_{33}-1
\end{array}\right\|}  \notag\\
& -\left[R_{23} R_{31}\left(R_{21}+1\right)\right]^{\frac{1}{2}}\left\|\begin{array}{lll}
R_{11}-1 & R_{12} & R_{13} \\
R_{21}+1 & R_{22}-1 & R_{23}-1 \\
R_{31}-1 & R_{32} & R_{33}
\end{array}\right\|=0 . \label{eq:8:6:37}\\
& \left(R_{11}+R_{12}+1\right)\left[R_{22}\right]^{\frac{1}{2}}\left\|\begin{array}{lll}
R_{11} & R_{12} & R_{13} \\
R_{21} & R_{22} & R_{23} \\
R_{31} & R_{32} & R_{33}
\end{array}\right\|+\left[R_{13}\left(R_{23}+1\right)\left(R_{12}+1\right)\right]^{\frac{1}{2}}\left\|\begin{array}{lll}
R_{11} & R_{12}+1 & R_{13}-1 \\
R_{21} & R_{22}-1 & R_{23}+1 \\
R_{31} & R_{32} & R_{33}
\end{array}\right\|  \notag\\
& -\left[R_{11} R_{33}(J+1)\right]^{\frac{1}{2}}\left\|\begin{array}{lll}
R_{11}-1 & R_{12} & R_{13} \\
R_{21} & R_{22}-1 & R_{23} \\
R_{31} & R_{32} & R_{33}-1
\end{array}\right\|=0 . \label{eq:8:6:38}\\
& \left(R_{22}-R_{33}\right)\left[R_{11} R_{23} R_{32}\right]^{\frac{1}{2}}\left\|\begin{array}{lll}
R_{11}-1 & R_{12} & R_{13} \\
R_{21} & R_{22} & R_{23}-1 \\
R_{31} & R_{32}-1 & R_{33}
\end{array}\right\|  \notag\\
& +\left(R_{33}-R_{11}\right)\left[R_{13} R_{22} R_{31}\right]^{\frac{1}{2}}\left\|\begin{array}{lll}
R_{11} & R_{12} & R_{13}-1 \\
R_{21} & R_{22}-1 & R_{23} \\
R_{31}-1 & R_{32} & R_{33}
\end{array}\right\|  \notag\\
& +\left(R_{11}-R_{22}\right)\left[R_{12} R_{21} R_{33}\right]^{\frac{1}{2}}\left\|\begin{array}{lll}
R_{11} & R_{12}-1 & R_{13} \\
R_{21}-1 & R_{22} & R_{23} \\
R_{31} & R_{32} & R_{33}-1
\end{array}\right\|=0 . \label{eq:8:6:39}
\end{align}
\section{SUMS OF PRODUCTS OF THE CLEBSCH-GORDAN COEFFICIENTS}\label{sec:8.7}\index{Clebsch-Gordan coefficient!sums of products of}\index{sums!of the Clebsch-Gordan coefficients}
Sums of products of the vector-addition coefficients are more conveniently represented in terms of the $3 j m$ symbols than in terms of the Clebsch-Gordan coefficients, because the $3 j m$ symbols have simpler symmetry relations. The sums expressed in terms of the $3 j m$ symbols will be given in Sec.~\ref{sec:12.2}. Nevertheless the ClebschGordan coefficients are also widely used. That is why we present the formulas which seem to be most useful for sums of products containing the Clebsch-Gordan coefficients. For practical convenience we consider various versions of the formulas which differ by permutations of upper and lower arguments. In these formulas we use the notation\index{dimension factor}\isymg{Pi-dim@$\Pi_{a b c \ldots}$, dimension factor $[(2a+1)(2b+1)\cdots]^{1/2}$}
\[
\Fact{a b \ldots c}=[(2 a+1)(2 b+1) \ldots(2 c+1)]^{\frac{1}{2}}
\]
\subsection{Sums Involving One Clebsch-Gordan Coefficient}\label{sec:8.7.1}\index{sums!with one Clebsch-Gordan coefficient}
\begin{gather}
\sum_{\alpha} \clebsch{a}{\alpha}{b}{0}{a}{\alpha}=\Fact{a}^{2} \delta_{b 0}, \label{eq:8:7:1}\\
\sum_{\alpha}(-1)^{a-\alpha} \clebsch{a}{\alpha}{a}{-\alpha}{c}{0}=\Fact{a} \delta_{c 0}, \label{eq:8:7:2}\\
\sum_{\alpha \beta \gamma}(-1)^{c+\gamma} \clebsch{a}{\alpha}{b}{\beta}{c}{\gamma}[(a+\alpha)!(a-\alpha)!(b+\beta)!(b-\beta)!(c+\gamma)!(c-\gamma)!]^{-\frac{1}{2}}=0. \label{eq:8:7:3}
\end{gather}
\subsection{Sums Involving Products of Two Clebsch-Gordan Coefficients}\label{sec:8.7.2}\index{sums!with two Clebsch-Gordan coefficients}
\begin{gather}
\sum_{\alpha \beta} \clebsch{a}{\alpha}{b}{\beta}{c}{\gamma} \clebsch{a}{\alpha}{b}{\beta}{c^{\prime}}{\gamma^{\prime}}=\delta_{c c^{\prime}} \delta_{\gamma \gamma^{\prime}}, \label{eq:8:7:4}\\
\sum_{\alpha \gamma} \clebsch{a}{\alpha}{b}{\beta}{c}{\gamma} \clebsch{a}{\alpha}{b^{\prime}}{\beta^{\prime}}{c}{\gamma}=\frac{\Fact{c}^{2}}{\Fact{b}^{2}} \delta_{b b^{\prime}} \delta_{\beta \beta^{\prime}}, \label{eq:8:7:5}\\
\sum_{\alpha \beta}(-1)^{b+\beta} \clebsch{a}{\alpha}{b}{\beta}{c}{\gamma} \clebsch{c^{\prime}}{-\gamma^{\prime}}{a}{\alpha}{b}{-\beta}=(-1)^{c+\gamma} \frac{\Fact{b}}{\Fact{c}} \delta_{c c^{\prime}} \delta_{\gamma \gamma^{\prime}}, \label{eq:8:7:6}\\
\sum_{\alpha \beta}(-1)^{a+\alpha} \clebsch{b}{\beta}{a}{\alpha}{c}{\gamma} \clebsch{a}{-\alpha}{c^{\prime}}{\gamma^{\prime}}{b}{\beta}=\frac{\Fact{b}}{\Fact{c}} \delta_{c c^{\prime}} \delta_{\gamma \gamma^{\prime}}, \label{eq:8:7:7}\\
\sum_{\alpha \beta} \clebsch{c}{\gamma}{b}{\beta}{a}{\alpha} \clebsch{c^{\prime}}{-\gamma^{\prime}}{a}{\alpha}{b}{\beta}=(-1)^{b-a-\gamma} \frac{\Fact{a b}}{\Fact{c}^{2}} \delta_{c c^{\prime}} \delta_{\gamma \gamma^{\prime}}, \label{eq:8:7:8}\\
\sum_{c \gamma} \clebsch{a}{\alpha}{b}{\beta}{c}{\gamma} \clebsch{a}{\alpha^{\prime}}{b}{\beta^{\prime}}{c}{\gamma}=\delta_{\alpha \alpha^{\prime}} \delta_{\beta \beta^{\prime}}, \label{eq:8:7:9}\\
\sum_{c \gamma} \Fact{c}^{2} \clebsch{a}{\alpha}{c}{\gamma}{b}{\beta} \clebsch{a}{\alpha^{\prime}}{c}{\gamma}{b}{\beta^{\prime}}=\Fact{b}^{2} \delta_{\alpha \alpha^{\prime}} \delta_{\beta \beta^{\prime}}, \label{eq:8:7:10}\\
\sum_{a \alpha}(-1)^{a-\alpha} \Fact{a}^{2} \clebsch{a}{\alpha}{b}{\beta}{c}{\gamma} \clebsch{a}{\alpha}{c}{\gamma^{\prime}}{b}{\beta^{\prime}}=\Fact{b c} \delta_{\gamma,-\gamma^{\prime}} \delta_{\beta,-\beta^{\prime}}. \label{eq:8:7:11}
\end{gather}
\subsection{Sums Involving Products of Three Clebsch-Gordan Coefficients}\label{sec:8.7.3}\index{sums!with three Clebsch-Gordan coefficients}
\begin{gather}
\sum_{\alpha \beta \delta} \clebsch{a}{\alpha}{b}{\beta}{c}{\gamma} \clebsch{d}{\delta}{b}{\beta}{e}{\varepsilon} \clebsch{a}{\alpha}{f}{\varphi}{d}{\delta}=\kappa_{1} \Fact{c d} \clebsch{c}{\gamma}{f}{\varphi}{e}{\varepsilon}\sixj{a}{b}{c}{e}{f}{d}, \label{eq:8:7:12}\\
\sum_{\alpha \beta \delta} \clebsch{b}{\beta}{c}{\gamma}{a}{\alpha} \clebsch{b}{\beta}{e}{\varepsilon}{d}{\delta} \clebsch{a}{\alpha}{f}{\varphi}{d}{\delta}=\kappa_{1} \frac{\Fact{a d d}}{\Fact{e}} \clebsch{c}{\gamma}{f}{\varphi}{e}{\varepsilon}\sixj{a}{b}{c}{e}{f}{d}, \label{eq:8:7:13}\\
\sum_{\alpha \beta \delta} 
\clebsch{b}{\beta}{a}{\alpha}{c}{\gamma} 
\clebsch{b}{\beta}{d}{\delta}{e}{\varepsilon} 
\clebsch{a}{\alpha}{f}{\varphi}{d}{\delta}
=\kappa_{2} \Fact{c d} 
\clebsch{c}{\gamma}{f}{\varphi}{e}{\varepsilon}
\sixj{a}{b}{c}{e}{f}{d}, \label{eq:8:7:14}\\
\sum_{\alpha \beta \delta}(-1)^{a-\alpha} \clebsch{a}{\alpha}{b}{\beta}{c}{\gamma} \clebsch{d}{\delta}{b}{\beta}{e}{\varepsilon} \clebsch{d}{\delta}{a}{-\alpha}{f}{\varphi}=\kappa_{1} \Fact{c f} \clebsch{c}{\gamma}{f}{\varphi}{e}{\varepsilon}\sixj{a}{b}{c}{e}{f}{d}, \label{eq:8:7:15}\\
\sum_{\alpha \beta \delta}(-1)^{b+\beta} \clebsch{a}{\alpha}{b}{\beta}{c}{\gamma} \clebsch{b}{-\beta}{e}{\varepsilon}{d}{\delta} \clebsch{a}{\alpha}{f}{\varphi}{d}{\delta}=\kappa_{1} \frac{\Fact{c d d}}{\Fact{e}} \clebsch{c}{\gamma}{f}{\varphi}{e}{\varepsilon}\sixj{a}{b}{c}{e}{f}{d}, \label{eq:8:7:16}\\
\sum_{\alpha \beta \delta}(-1)^{a-\alpha} \clebsch{b}{\beta}{a}{\alpha}{c}{\gamma} \clebsch{b}{\beta}{d}{\delta}{e}{\varepsilon} \clebsch{d}{\delta}{a}{-\alpha}{f}{\varphi}=\kappa_{2} \Fact{c f} \clebsch{c}{\gamma}{f}{\varphi}{e}{\varepsilon}
\sixj{a}{b}{c}{e}{f}{d}, \label{eq:8:7:17}\\
\sum_{\alpha \beta \delta}(-1)^{b+\beta} \clebsch{b}{-\beta}{c}{\gamma}{a}{\alpha} \clebsch{d}{\delta}{b}{\beta}{e}{\varepsilon} \clebsch{a}{\alpha}{f}{\varphi}{d}{\delta}=\kappa_{1} \Fact{a d} \clebsch{c}{\gamma}{f}{\varphi}{e}{\varepsilon}
\sixj{a}{b}{c}{e}{f}{d}, \label{eq:8:7:18}\\
\sum_{\alpha \beta \delta}(-1)^{a-\alpha} \clebsch{a}{\alpha}{c}{-\gamma}{b}{\beta} \clebsch{d}{\delta}{b}{-\beta}{e}{\varepsilon} \clebsch{a}{\alpha}{f}{\varphi}{d}{\delta}=\kappa_{1} \Fact{b d} \clebsch{c}{\gamma}{f}{\varphi}{e}{\varepsilon}\sixj{a}{b}{c}{e}{f}{d}, \label{eq:8:7:19}
\end{gather}
where $\kappa_{1}=(-1)^{b+c+d+f}, \kappa_{2}=(-1)^{a+b+e+f}$.\\
\subsection{Sums Involving Products of Four Clebsch-Gordan Coefficients}\label{sec:8.7.4}\index{sums!with four Clebsch-Gordan coefficients}
\begin{align}
& \sum_{\beta \gamma \varepsilon \varphi} \clebsch{b}{\beta}{c}{\gamma}{a}{\alpha} \clebsch{e}{\varepsilon}{f}{\varphi}{d}{\delta} \clebsch{e}{\varepsilon}{b}{\beta}{g}{\eta} \clebsch{f}{\varphi}{c}{\gamma}{j}{\mu}  \notag\\
& =(-1)^{a-b+c+d+e-f} \sum_{s \sigma} \Fact{s s a g} \clebsch{a}{\alpha}{s}{\sigma}{j}{\mu} \clebsch{g}{\eta}{s}{\sigma}{d}{\delta}\sixj{b}{c}{a}{j}{s}{f}\sixj{b}{e}{g}{d}{s}{f}=\Fact{a d g j} \sum_{k \kappa} \clebsch{g}{\eta}{j}{\mu}{k}{\kappa} \clebsch{d}{\delta}{a}{\alpha}{k}{\kappa}\ninej{c}{b}{a}{f}{e}{d}{j}{g}{k}, \label{eq:8:7:20}\\
& \sum_{\beta \gamma \varepsilon \varphi} \clebsch{b}{\beta}{c}{-\gamma}{a}{\alpha} \clebsch{e}{\varepsilon}{f}{-\varphi}{d}{\delta} \clebsch{b}{\beta}{e}{\varepsilon}{g}{\eta} \clebsch{c}{\gamma}{f}{\varphi}{j}{\mu}=(-1)^{b+e-g} \Fact{a d g j} \sum_{k \kappa} \clebsch{g}{\eta}{j}{-\mu}{k}{\kappa} \clebsch{d}{\delta}{a}{\alpha}{k}{\kappa}\ninej{c}{b}{a}{f}{e}{d}{j}{g}{k}, \label{eq:8:7:21}\\
& \sum_{\beta \gamma \varepsilon \varphi} \clebsch{b}{\beta}{a}{\alpha}{c}{\gamma} \clebsch{f}{\varphi}{j}{\mu}{c}{\gamma} \clebsch{b}{\beta}{g}{\eta}{e}{\varepsilon} \clebsch{f}{\varphi}{d}{\delta}{e}{\varepsilon}=(-1)^{a-b+f-j} \Fact{c e}^{2} \sum_{k \kappa} \clebsch{g}{\eta}{j}{\mu}{k}{\kappa} \clebsch{d}{\delta}{a}{\alpha}{k}{\kappa}\ninej{c}{b}{a}{f}{e}{d}{j}{g}{k}, \label{eq:8:7:22}\\
& \sum_{\beta \gamma \varepsilon \varphi} \clebsch{a}{\alpha}{b}{\beta}{c}{\gamma} \clebsch{g}{\eta}{e}{\varepsilon}{b}{\beta} \clebsch{d}{\delta}{f}{\varphi}{e}{\varepsilon} \clebsch{j}{\mu}{c}{\gamma}{f}{\varphi}=(-1)^{d+e-c-j} \Fact{b c e f} \sum_{k \kappa}(-1)^{k-\kappa} \clebsch{g}{\eta}{j}{\mu}{k}{-\kappa} \clebsch{d}{\delta}{a}{\alpha}{k}{\kappa}\ninej{c}{b}{a}{f}{e}{d}{j}{g}{k}, \label{eq:8:7:23}\\
& \sum_{\beta \gamma \varepsilon \varphi} \clebsch{b}{\beta}{a}{-\alpha}{c}{\gamma} \clebsch{e}{\varepsilon}{d}{-\delta}{f}{\varphi} \clebsch{g}{\eta}{b}{-\beta}{e}{\varepsilon} \clebsch{j}{\mu}{c}{-\gamma}{f}{\varphi}=(-1)^{b-c-g-\alpha+\eta} \Fact{c e f f} \sum_{k \kappa} \clebsch{g}{\eta}{j}{-\mu}{k}{\kappa} \clebsch{d}{\delta}{a}{\alpha}{k}{\kappa}\ninej{c}{b}{a}{f}{e}{d}{j}{g}{k}, \label{eq:8:7:24}
\end{align}
\begin{align}
& \sum_{\beta \gamma \varepsilon \varphi} \clebsch{b}{\beta}{c}{-\gamma}{a}{\alpha} \clebsch{e}{\varepsilon}{f}{\varphi}{d}{\delta} \clebsch{e}{-\varepsilon}{g}{\eta}{b}{\beta} \clebsch{f}{\varphi}{j}{\mu}{c}{\gamma}=(-1)^{b+f-g-\delta} \Fact{a b c d} \sum_{k \kappa} \clebsch{g}{\eta}{j}{-\mu}{k}{\kappa} \clebsch{d}{\delta}{a}{\alpha}{k}{\kappa}\ninej{c}{b}{a}{f}{e}{d}{j}{g}{k}, \label{eq:8:7:25}\\
& \sum_{\beta \gamma \varepsilon \varphi} \clebsch{b}{\beta}{c}{\gamma}{a}{\alpha} \clebsch{e}{\varepsilon}{f}{\varphi}{d}{\delta} \clebsch{e}{\varepsilon}{g}{\eta}{b}{\beta} \clebsch{f}{\varphi}{j}{\mu}{c}{\gamma}=\sum_{k \kappa} \Fact{b c d k} \clebsch{g}{\eta}{j}{\mu}{k}{\kappa} \clebsch{d}{ \delta}{k}{\kappa}{a}{\alpha}\ninej{a}{b}{c}{d}{e}{f}{k}{g}{j}, \label{eq:8:7:26}\\
& \sum_{\beta \gamma \varepsilon \varphi} \clebsch{b}{\beta}{c}{-\gamma}{a}{\alpha} \clebsch{e}{\varepsilon}{f}{-\varphi}{d}{\delta} \clebsch{g}{\eta}{b}{-\beta}{e}{\varepsilon} \clebsch{j}{\mu}{f}{-\varphi}{c}{\gamma}=(-1)^{c+e-g+j+\alpha-\mu} \Fact{a d e c} \sum_{k \kappa} \clebsch{g}{\eta}{j}{-\mu}{k}{\kappa} \clebsch{d}{\delta}{a}{\alpha}{k}{\kappa}\ninej{c}{b}{a}{f}{e}{d}{j}{g}{k}, \label{eq:8:7:27}\\
& \sum_{\beta \gamma \varepsilon \varphi} \clebsch{b}{\beta}{c}{\gamma}{a}{\alpha} \clebsch{b}{\beta}{g}{\eta}{e}{\varepsilon} \clebsch{f}{\varphi}{d}{\delta}{e}{\varepsilon} \clebsch{f}{\varphi}{c}{\gamma}{j}{\mu}=(-1)^{j-a+\delta-\eta} \Fact{a e e j} \sum_{k \kappa} \clebsch{g}{\eta}{j}{-\mu}{k}{\kappa} \clebsch{d}{\delta}{a}{-\alpha}{k}{\kappa}\ninej{c}{b}{a}{f}{e}{d}{j}{g}{k}, \label{eq:8:7:28}\\
& \sum_{\beta \gamma \varepsilon \varphi} \clebsch{b}{\beta}{c}{\gamma}{a}{\alpha} \clebsch{g}{\eta}{e}{\varepsilon}{b}{\beta} \clebsch{f}{\varphi}{d}{\delta}{e}{\varepsilon} \clebsch{f}{\varphi}{c}{\gamma}{j}{\mu}=(-1)^{j-a+g+\delta} \Fact{a b e j} \sum_{k \kappa} \clebsch{g}{-\eta}{j}{-\mu}{k}{\kappa} \clebsch{d}{\delta}{a}{-\alpha}{k}{\kappa}\ninej{c}{b}{a}{f}{e}{d}{j}{g}{k}, \label{eq:8:7:29}\\
& \sum_{\beta \gamma \varepsilon\varphi}(-1)^{c-\gamma+e-\varepsilon} \clebsch{a}{\alpha}{b}{\beta}{c}{\gamma} \clebsch{d}{\delta}{f}{\varphi}{e}{\varepsilon} \clebsch{e}{\varepsilon}{b}{\beta}{g}{\eta} \clebsch{c}{\gamma}{f}{\varphi}{j}{\mu}=(-1)^{a+d-\alpha-\delta} \Fact{c e g j} \sum_{k \kappa} \clebsch{g}{\eta}{j}{-\mu}{k}{\kappa} \clebsch{d}{\delta}{a}{-\alpha}{k}{\kappa}\ninej{c}{b}{a}{f}{e}{d}{j}{g}{k} . \label{eq:8:7:30}
\end{align}
\subsection{Sums Involving Products of the Clebsch-Gordan Coefficients and One 6j Symbol}\label{sec:8.7.5}\index{sums!with a 6j symbol@with a $6j$ symbol}\index{6j symbol@$6j$ symbol}
\begin{align}
& \sum_{e \varepsilon}(-1)^{2 e} \Fact{c d} \clebsch{b}{\beta}{d}{\delta}{e}{\varepsilon} \clebsch{f}{\varphi}{c}{\gamma}{e}{\varepsilon}
\sixj{a}{b}{c}{e}{f}{d}
=\clebsch{a}{\alpha}{b}{\beta}{c}{\gamma} 
\clebsch{a}{\alpha}{f}{\varphi}{d}{\delta}, \label{eq:8:7:31}\\
& \sum_{f \varphi}(-1)^{c+d+f} \Fact{c e} \clebsch{e}{\varepsilon}{a}{\alpha}{f}{\varphi} \clebsch{d}{\delta}{c}{\gamma}{f}{\varphi}\sixj{b}{a}{c}{f}{d}{e}=\clebsch{a}{\alpha}{b}{\beta}{c}{\gamma} \clebsch{d}{\delta}{b}{\beta}{e}{\varepsilon}, \label{eq:8:7:32}\\
& \sum_{c \gamma}(-1)^{2 e-d+\alpha+\varphi} \Fact{a e} \clebsch{f}{-\varphi}{b}{\beta}{c}{\gamma} \clebsch{e}{\varepsilon}{a}{-\alpha}{c}{\gamma}\sixj{c}{f}{b}{d}{e}{a}=\clebsch{b}{\beta}{d}{\delta}{e}{\varepsilon} \clebsch{f}{\varphi}{d}{\delta}{a}{\alpha}, \label{eq:8:7:33}\\
& \sum_{c \gamma}(-1)^{c+d-\beta-\varphi} \Fact{d}^{2} \clebsch{a}{\alpha}{b}{\beta}{c}{\gamma} \clebsch{f}{-\varphi}{e}{\varepsilon}{c}{\gamma}\sixj{a}{b}{c}{e}{f}{d}=\clebsch{a}{\alpha}{f}{\varphi}{d}{\delta} \clebsch{b}{-\beta}{e}{\varepsilon}{d}{\delta}, \label{eq:8:7:34}\\
& \sum_{c \gamma}(-1)^{2 e} \Fact{c d} \clebsch{a}{\alpha}{b}{\beta}{c}{\gamma} \clebsch{f}{\varphi}{c}{\gamma}{e}{\varepsilon}\sixj{a}{b}{c}{e}{f}{d}=\clebsch{b}{\beta}{d}{\delta}{e}{\varepsilon} \clebsch{a}{\alpha}{f}{\varphi}{d}{\delta}, \label{eq:8:7:35}\\
& \sum_{f \varphi}(-1)^{2 c} \Fact{e f} \clebsch{b}{\beta}{d}{\delta}{f}{\varphi} \clebsch{a}{\alpha}{f}{\varphi}{c}{\gamma}\sixj{a}{b}{e}{d}{c}{f}=\clebsch{b}{\beta}{a}{\alpha}{e}{\varepsilon} \clebsch{d}{\delta}{e}{\varepsilon}{c}{\gamma}, \label{eq:8:7:36}\\
& \sum_{c \gamma}(-1)^{e+d-\beta} \frac{\Fact{c d d}}{\Fact{e}} \clebsch{a}{\alpha}{b}{\beta}{c}{\gamma} \clebsch{f}{\varphi}{c}{\gamma}{e}{\varepsilon}\sixj{a}{b}{c}{e}{f}{d}=\clebsch{b}{-\beta}{e}{\varepsilon}{d}{\delta} \clebsch{a}{\alpha}{f}{\varphi}{d}{\delta}, \label{eq:8:7:37}\\
& \sum_{c \gamma}(-1)^{2 e} \frac{\Fact{a c c}}{\Fact{b}} \clebsch{f}{\varphi}{c}{\gamma}{b}{\beta} \clebsch{c}{\gamma}{a}{\alpha}{e}{\varepsilon}\sixj{c}{f}{b}{d}{e}{a}=\clebsch{d}{\delta}{b}{\beta}{e}{\varepsilon} \clebsch{f}{\varphi}{d}{\delta}{a}{\alpha} . \label{eq:8:7:38}
\end{align}
\subsection{Sums Involving Products of the Clebsch-Gordan Coefficients and One 9j Symbol}\label{sec:8.7.6}\index{sums!with a 9j symbol@with a $9j$ symbol}\index{9j symbol@$9j$ symbol}
\begin{align}
& \sum_{a k} \Fact{a d g j} \clebsch{b}{\beta}{c}{\gamma}{a}{\alpha} \clebsch{g}{\eta}{j}{\mu}{k}{\kappa} \clebsch{d}{\delta}{a}{\alpha}{k}{\kappa}\ninej{c}{b}{a}{f}{e}{d}{j}{g}{k}=\clebsch{e}{\varepsilon}{f}{\varphi}{d}{\delta} \clebsch{e}{\varepsilon}{b}{\beta}{g}{\eta} \clebsch{f}{\varphi}{c}{\gamma}{j}{\mu}, \label{eq:8:7:39}\\
& \sum_{g j} \Fact{a d g j} \clebsch{g}{\eta}{j}{\mu}{k}{\kappa} \clebsch{e}{\varepsilon}{b}{\beta}{g}{\eta} \clebsch{f}{\varphi}{c}{\gamma}{j}{\mu}\ninej{c}{b}{a}{f}{e}{d}{j}{g}{k}=\clebsch{d}{\delta}{a}{\alpha}{k}{\kappa} \clebsch{e}{\varepsilon}{f}{\varphi}{d}{\delta} \clebsch{b}{\beta}{c}{\gamma}{a}{\alpha}, \label{eq:8:7:40}\\
& \sum_{g j}(-1)^{b+e-g} \Fact{a d g j} \clebsch{g}{\eta}{j}{-\mu}{k}{\kappa} \clebsch{b}{\beta}{e}{\varepsilon}{g}{\eta} \clebsch{c}{\gamma}{f}{\varphi}{j}{\mu}\ninej{c}{b}{a}{f}{e}{d}{j}{g}{k}=\clebsch{d}{\delta}{a}{\alpha}{k}{\kappa} \clebsch{e}{\varepsilon}{f}{-\varphi}{d}{\delta} \clebsch{b}{\beta}{c}{-\gamma}{a}{\alpha}, \label{eq:8:7:41}\\
& \sum_{a d}(-1)^{b+f-g-\delta} \Fact{a b c d} \clebsch{b}{\beta}{c}{-\gamma}{a}{\alpha} \clebsch{e}{\varepsilon}{f}{\varphi}{d}{\delta} \clebsch{d}{\delta}{a}{\alpha}{k}{\kappa}\ninej{c}{b}{a}{f}{e}{d}{j}{g}{k}=\clebsch{g}{\eta}{j}{-\mu}{k}{\kappa} \clebsch{e}{-\varepsilon}{g}{\eta}{b}{\beta} \clebsch{f}{\varphi}{j}{\mu}{c}{\gamma}, \label{eq:8:7:42}\\
& \sum_{g j}(-1)^{g-b-f+\varphi+\varepsilon} \frac{\Fact{a d j j g g}}{\Fact{b c}} \clebsch{g}{\eta}{j}{-\mu}{k}{\kappa} \clebsch{e}{-\varepsilon}{g}{\eta}{b}{\beta} \clebsch{f}{\varphi}{j}{\mu}{c}{\gamma}\ninej{c}{b}{a}{f}{e}{d}{j}{g}{k}=\clebsch{b}{\beta}{c}{-\gamma}{a}{\alpha} \clebsch{e}{\varepsilon}{f}{\varphi}{d}{\delta} \clebsch{d}{\delta}{a}{\alpha}{k}{\kappa} . \label{eq:8:7:43}
\end{align}
Sums of larger numbers of the Clebsch-Gordan coefficients may be obtained from the above equations using the orthogonality relation (\eqref{eq:8:1:8}).

\subsection{Some Additional Sums of Products of Two Clebsch-Gordan Coefficients}\label{sec:8.7.7}\index{sums!more, with two Clebsch-Gordan coefficients}
We note also some other sums of products of the Clebsch-Gordan coefficients. The first two sums, obtained by Morgan \cite{ref136}, are written as
\begin{equation}
\sum_{l^{\prime}=0}^{l} \frac{1}{\left(2 l^{\prime}-1\right)\left(2 l-2 l^{\prime}+2 J+1\right)}\left(\clebsch{l}{0}{l^{\prime}+J}{0}{l-l^{\prime}+J}{0}\right)^{2}=-\frac{\delta_{l0}}{2 J+1}, \label{eq:8:7:44}
\end{equation}
and
\begin{equation}
\sum_{l^{\prime}=0}^{l}\left(\frac{1}{2 l^{\prime}+3}-\frac{l+1}{2 l+3} \frac{1}{2 l^{\prime}+1}\right) \frac{1}{2 l-2 l^\prime+2 J+1}
\left(
  \clebsch{l}{0}{l^\prime+J}{0}{l-l^\prime+J}{0}
\right)^{2}=0 . \label{eq:8:7:45}
\end{equation}
In these equations $l, l^{\prime}, J$ are integers.\\[0pt]
Din proved the following identity \cite{ref134}
\begin{equation}
\sum_{i=|c-b|}^{c+b} \frac{2 i+1}{i(i+1)-a(a+1)}\left(\clebsch{i}{0}{b}{0}{c}{0}\right)^{2}=0, \label{eq:8:7:46}
\end{equation}
where $a, b, c$ are non-negative integers, with $a+b+c$ odd, which satisfy the triangular condition
\[
|c-b| \leq a \leq c+b .
\]
The sum runs only over values of $i$ for which the Clebsch-Gordan coefficient does not vanish. Dunlop and Judd \cite{ref135} give the following relation
\begin{equation}
\sum_{m} \clebsch{a}{m}{k}{0}{a}{m} \clebsch{c}{M-m}{k}{0}{c}{M-m}=(-1)^{k} \frac{\Fact{a c}}{2 k+1} \sqrt{\frac{(2 a-k)!(2 c+k+1)!}{(2 c-k)!(2 a+k+1)!}}, \label{eq:8:7:47}
\end{equation}
which is valid for $a-c \geq|M|$.

\section{GENERATING FUNCTIONS}\label{sec:8.8}\index{Clebsch-Gordan coefficient!generating functions}\index{generating functions!for the Clebsch-Gordan coefficients}
The Clebsch-Gordan coefficients may be generated by expansions of some functions in power series or finite sums. Below we present several generating functions. In these relations $J \equiv a+b+c$.

\subsection{Regge Determinant to Power $J$}\label{sec:8.8.1}\index{Regge determinant}\index{generating functions!Regge determinant}
\begin{equation}
\left|\begin{array}{ccc}
u_{1} & u_{2} & u_{3}  \\
v_{1} & v_{2} & v_{3} \\
w_{1} & w_{2} & w_{3}
\end{array}\right|^{J}=J!\sqrt{(J+1)!} \sum_{R_{i k}}\left\|\begin{array}{lll}
R_{11} & R_{12} & R_{13} \\
R_{21} & R_{22} & R_{23} \\
R_{31} & R_{32} & R_{33}
\end{array}\right\| \frac{u_{1}^{R_{11}} u_{2}^{R_{12}} u_{3}^{R_{13}} v_{1}^{R_{21}} v_{2}^{R_{22}} v_{3}^{R_{23}} w_{1}^{R_{31}} w_{2}^{R_{32}} w_{3}^{R_{33}}}{\left[R_{11}!R_{12}!R_{13}!R_{21}!R_{22}!R_{23}!R_{31}!R_{32}!R_{33}!\right]^{\frac{1}{2}}} . \label{eq:8:8:1}
\end{equation}
(Regge \cite{ref094}, Shelepin \cite{ref106}). Here $R_{i k}$ are the elements of the $R$-symbol (see Sec.~\ref{sec:8.1.3}). The $R$-symbols and the Clebsch-Gordan coefficients are related by Eqs.~\eqref{eq:8:1:14} and \eqref{eq:8:1:11}.

\subsection{Products of Binomials}\label{sec:8.8.2}\index{generating functions!products of binomials}
\begin{align}
\MoveEqLeft{(t_1-t_2)^{J-2c} (t_2-t_3)^{J-2a}(t_3-t_1)^{J-2b}}\notag\\
= & \left[(J+1)!\frac{(J-2 a)!(J-2 b)!(J-2 c)!}{2 c+1}\right]^{\frac{1}{2}} 
\sum_{\substack{\alpha, \beta, \gamma \\ \alpha+\beta+\gamma=0}}(-1)^{a-b-\gamma}   %\notag\\\times & 
\frac{t_{1}^{a+\alpha} t_{2}^{b+\beta} t_{3}^{c+\gamma}}{[(a+\alpha)!(a-\alpha)!(b+\beta)!(b-\beta)!(c+\gamma)!(c-\gamma)!]^{\frac{1}{2}}} \clebsch{a}{\alpha}{b}{\beta}{c}{-\gamma}, \label{eq:8:8:2}\\
\MoveEqLeft{\left(v_{1} u_{2}-u_{1} v_{2}\right)^{J-2 c}\left(w_{1} v_{2}-w_{2} v_{1}\right)^{J-2 a}\left(u_{1} w_{2}-w_{1} u_{2}\right)^{J-2 b}} \notag \\
= & \left[\frac{(J+1)!(J-2 a)!(J-2 b)!(J-2 c)!}{2 c+1}\right]^{\frac{1}{2}}\notag\\& \times 
\sum_{\substack{\alpha, \beta, \gamma \\ \alpha+\beta+\gamma=0}}
(-1)^{a-b-\gamma} 
\frac{u_{1}^{a-\alpha} u_{2}^{a+\alpha} v_{1}^{b-\beta} v_{2}^{b+\beta} w_{1}^{c-\gamma} w_{2}^{c+\gamma}}{[(a+\alpha)!(a-\alpha)!(b+\beta)!(b-\beta)!(c+\gamma)!(c-\gamma)!]^{\frac{1}{2}}} \clebsch{a}{\alpha}{b}{\beta}{c}{-\gamma} . \label{eq:8:8:3}
\end{align}

\subsection{Exponential Function}\label{sec:8.8.3}\index{generating functions!exponential function}
\begin{align}
\MoveEqLeft{\exp\{x_1(t_2-t_3)+x_2(t_3-t_1)+x_3(t_1-t_2)\}}\notag\\
  &=\sum_{\substack{a, b, c \\ \alpha, \beta, \gamma}} \frac{(-1)^{a-b-\gamma}}{\Delta(a b c)} \frac{x_1^{-a+b+c}x_2^{a-b+c}x_3^{a+b-c}t_1^{a+\alpha}t_2^{b+\beta}t_3^{c+\gamma}}{[(2 c+1)(a+\alpha)!(a-\alpha)!(b+\beta)!(b-\beta)!(c+\gamma)!(c-\gamma)!]^{\frac{1}{2}}} \clebsch{a}{\alpha}{b}{\beta}{c}{-\gamma}, \label{eq:8:8:4}
\end{align}
where $\Delta(a b c)$ is given by \eqref{eq:8:2:1}.

\subsection{Hypergeometric Function}\label{sec:8.8.4}\index{generating functions!hypergeometric function}
\begin{gather}
{\left[\frac{(J-2 b)!(c+\gamma)!(2 c+1)}{(J-2 a)!(J-2 c)!(c-\gamma)!(J+1)!}\right]^{\frac{1}{2}} \frac{(t-1)^{J-2 c}}{(a-b+\gamma)!} F(\gamma-c, a-b-c ; a-b+\gamma+1 ; t)}  \notag\\
=\sum_{\substack{\alpha, \beta \\
\alpha+\beta=\gamma}} \clebsch{a}{\alpha}{b}{\beta}{c}{\gamma} \frac{t^{b-\beta}}{[(a+\alpha)!(a-\alpha)!(b+\beta)!(b-\beta)!]^{\frac{1}{2}}}. \label{eq:8:8:5}
\end{gather}
(Akim and Levin \cite{ref046})

\subsection{Wigner D-Function}\label{sec:8.8.5}\index{generating functions!Wigner $D$-function}\isym{dJMM@$d_{M M^{\prime}}^{J}(\beta)$, Wigner (small) $d$-function}
\begin{equation}
\left[\frac{(2 c+1)}{(J+1)!(J-2 c)!}\right]^{\frac{1}{2}} d_{\gamma \gamma^{\prime}}^{c}(\theta)=\sum_{\substack{\alpha, \beta \\ \alpha+\beta=\gamma}}(-1)^{b+\beta} \clebsch{a}{\alpha}{b}{\beta}{c}{\gamma} \frac{\left(\cos \frac{\theta}{2}\right)^{a+b+\alpha-\beta}\left(\sin \frac{\theta}{2}\right)^{a+b-\alpha+\beta}}{[(a+\alpha)!(a-\alpha)!(b+\beta)!(b-\beta)!]^{\frac{1}{2}}}, \label{eq:8:8:6}
\end{equation}
where $a-b=\gamma^{\prime},|a-b| \leq c \leq a+b$.

\subsection{Generating Function for the Coefficients $\clebsch{a}{0}{b}{0}{c}{0}$}\label{sec:8.8.6}\index{generating functions!for coefficients with zero projections}
\begin{equation}
\frac{1}{1+u^{2}+v^{2}+w^{2}}=\sum_{a b c}(-1)^{a-b}\left[\frac{(J+1)!}{2 c+1}\right]^{\frac{1}{2}} \clebsch{a}{0}{b}{0}{c}{0} \frac{u^{J-2 a} v^{J-2 b} w^{J-2 c}}{[(J-2 a)!(J-2 b)!(J-2 c)!]^{\frac{1}{2}}} . \label{eq:8:8:7}
\end{equation}
(Schwinger \cite{ref101})

\section{CLASSICAL LIMIT AND ASYMPTOTIC EXPRESSIONS FOR THE CLEBSCH-GORDAN COEFFICIENTS}\label{sec:8.9}\index{Clebsch-Gordan coefficient!classical limit}\index{Clebsch-Gordan coefficient!asymptotic expressions}\index{asymptotics!of the Clebsch-Gordan coefficients}\index{classical limit}
Asymptotic expressions permit us to simplify calculations of the Clebsch-Gordan coefficients. Note that these expressions appear to be accurate even if the arguments are not very large.\\
\subsection{Asymptotic Expressions for $a, c \gg b$}\label{sec:8.9.1}\index{asymptotics!for two large momenta}
\begin{multline}
\clebsch{a}{\alpha}{b}{\beta}{c}{\gamma} \approx[(b+\beta)!(b-\beta)!(b+\delta)!(b-\delta)!]^{\frac{1}{2}}  \\
\times \sum_{s} \frac{(-1)^{b-\delta-s}\left(\cos \frac{\vartheta}{2}\right)^{2s+\beta+\delta}\left(\sin \frac{\vartheta}{2}\right)^{2 b-\beta-\delta-2 s}}{s!(s+\beta+\delta)!(b-\beta-s)!(b-\delta-s)!}=\delta_{\gamma-\alpha, \beta} D_{\beta \delta}^{b}(0, \vartheta, 0), \label{eq:8:9:1}
\end{multline}
where $\delta=c-a$, and $\vartheta$ is the angle between the momentum $c$ and the $z$ axis,
\begin{equation}
\cos \frac{\vartheta}{2}=\sqrt{\frac{c+\gamma+\frac{1}{2}}{2 c+1}}=\clebsch{c}{\gamma}{\frac{1}{2}}{\frac{1}{2}}{c+\frac{1}{2}}{\gamma+\frac{1}{2}}, \quad \sin \frac{\vartheta}{2}=\sqrt{\frac{c-\gamma+\frac{1}{2}}{2 c+1}}=\clebsch{c}{\gamma}{\frac{1}{2}}{-\frac{1}{2}}{c+\frac{1}{2}}{\gamma-\frac{1}{2}}, \quad \cos \vartheta=\frac{\gamma}{c+\frac{1}{2}} . \label{eq:8:9:2}
\end{equation}
This formula was obtained by Edmonds \cite{ref064}, Brussaard and Tolhoek \cite{ref060}. In particular,
\begin{align}
\clebsch{a}{\alpha}{b}{b}{c}{\gamma} &\approx \delta_{b, \gamma-\alpha}(-1)^{b-\delta}\left[\frac{(2 b)!}{(b+\delta) !(b-\delta)!}\right]^{\frac{1}{2}}\left(\cos \frac{\vartheta}{2}\right)^{b+\delta}\left(\sin \frac{\vartheta}{2}\right)^{b-\delta}, \label{eq:8:9:3}\\
\clebsch{a}{\alpha}{b}{\beta}{c}{c} &\approx \delta_{\beta, c-\alpha} D_{\beta c-a}^{b}(0,0,0)=\delta_{\beta, c-\alpha} \delta_{\beta, c-a}, \label{eq:8:9:4}\\
\clebsch{a}{-\beta}{b}{\beta}{c}{0} &\approx D_{\beta c-a}^{b}\left(0, \frac{\pi}{2}, 0\right), \label{eq:8:9:5}\\
\clebsch{c}{\gamma}{b}{0}{c}{\gamma} &\approx P_{b}(\cos \vartheta), \label{eq:8:9:6}\\
\clebsch{c}{\alpha}{b}{\beta}{c}{\gamma} &\approx \delta_{\beta, \gamma-\alpha} \sqrt{\frac{4 \pi}{2 b+1}} Y_{b \beta}(\vartheta, 0), \label{eq:8:9:7}\\
\clebsch{c}{\pm c}{b}{0}{c}{\pm c} &\approx( \pm 1)^{b}. \label{eq:8:9:8}
\end{align}
If $a, c \gg b \gg 1$ one can use in \eqref{eq:8:9:1} the asymptotic expressions for the Wigner $D$-function (Eqs.~\eqref{eq:4:18:1}\eqref{eq:4:18:4}). Then the Clebsch-Gordan coefficients may be represented in the form
\begin{gather}
\clebsch{a}{\alpha}{b}{\beta}{c}{\gamma} \approx \frac{(-1)^{\beta-\delta}}{(\beta-\delta)!}\left[\frac{2}{\pi} \frac{(b+\beta)!(b-\beta+1)!}{(b+\delta)!(b-\delta)!}\right]^{\frac{1}{2}} \frac{\cos \left[(2 b+1) \frac{\vartheta}{2}-\frac{\pi}{2}\left(\beta-\delta+\frac{1}{2}\right)\right]}{\sqrt{\sin \vartheta}}, \label{eq:8:9:9}\\
\text { if } \beta \geq \delta, \beta \geq 0, \notag\\
\clebsch{a}{\alpha}{b}{\beta}{c}{\gamma} \approx \frac{1}{(\delta-\beta)!}\left[\frac{2}{\pi} \frac{(b+\delta)!(b-\delta+1)!}{(b+\beta)!(b-\beta)!}\right]^{\frac{1}{2}} \frac{\cos \left[(2 b+1) \frac{\vartheta}{2}-\frac{\pi}{2}\left(\delta-\beta+\frac{1}{2}\right)\right]}{\sqrt{\sin \vartheta}}, \label{eq:8:9:10}\\
\text { if } \delta \geq \beta, \delta \geq 0 . \notag
\end{gather}
\subsection{Asymptotic Expressions for $a, b, c, \gamma \gg a+b-\gamma$}\label{sec:8.9.2}\index{asymptotics!for large momenta and projections}
\begin{equation}
\clebsch{a}{\alpha}{b}{\beta}{c}{\gamma} \approx(-1)^{a+b-c} \frac{(c+\gamma+1)^{b-\beta}}{(a-b-\alpha)!}\left[\frac{(2 c+1)(a-b+c)!(a-\alpha)!(b+\beta)!(c-\gamma)!(c+\gamma)!}{(a+b-c)!(-a+b+c)!(a+b+c+1)!(a+\alpha)!(b-\beta)!}\right]^{\frac{1}{2}}. \label{eq:8:9:11}
\end{equation}
(Akim and Levin \cite{ref046}).

\subsection{Semiclassical Formulas for $a, b, c \gg 1$}\label{sec:8.9.3}\index{semiclassical formulas}
\begin{equation}
\clebsch{a}{\alpha}{b}{\beta}{c}{\gamma} \approx(-1)^{2 b+c+\gamma+1} \sqrt{\frac{2 c+1}{2 \pi S}} \cos \left[j_{1} \theta_{1}+j_{2} \theta_{2}+j_{3} \theta_{3}-m_{2} \varphi_{1}+m_{1} \varphi_{2}+\frac{\pi}{4}\right], \label{eq:8:9:12}
\end{equation}
(Ponzano and Regge \cite{ref089}), where
\begin{gather}
j_{1} \equiv a+\frac{1}{2}, j_{2} \equiv b+\frac{1}{2}, j_{3} \equiv c+\frac{1}{2}, m_{1} \equiv \alpha, \quad m_{2} \equiv \beta, \quad m_{3} \equiv-\gamma, \label{eq:8:9:13}\\
S^{2}=\frac{1}{16}\left\{-j_{1}^{4}-j_{2}^{4}-j_{3}^{4}+2\left(j_{1}^{2} j_{2}^{2}+j_{1}^{2} j_{3}^{2}+j_{2}^{2} j_{3}^{2}\right)+4\left(j_{1}^{2} m_{2} m_{3}+j_{2}^{2} m_{1} m_{3}+j_{3}^{2} m_{1} m_{2}\right)\right\}  \notag\\
=-\frac{1}{16}\left|\begin{array}{cccc}
0 & j_{1}^{2}-m_{1}^{2} & j_{2}^{2}-m_{2}^{2} & 1 \\
j_{1}^{2}-m_{1}^{2} & 0 & j_{3}^{2}-m_{3}^{2} & 1 \\
j_{2}^{2}-m_{2}^{2} & j_{3}^{2}-m_{3}^{2} & 0 & 1 \\
1 & 1 & 1 & 0
\end{array}\right| . \label{eq:8:9:14}
\end{gather}
If $m_{1}=m_{2}=m_{3}=0$, \eqref{eq:8:9:14} reduces to
\begin{gather}
S^{2}=\frac{1}{16}\left(a+b+c+\frac{3}{2}\right)\left(-a+b+c+\frac{1}{2}\right)\left(a-b+c+\frac{1}{2}\right)\left(a+b-c+\frac{1}{2}\right), \label{eq:8:9:15}
\end{gather}
\begin{equation}\begin{gathered}
\cos \theta_{i}=\frac{2 j_{i}^{2} m_{l}+m_{i}\left(j_{i}^{2}-j_{k}^{2}+j_{l}^{2}\right)}{\sqrt{\left(j_{i}^{2}-m_{i}^{2}\right)\left(4 j_{i}^{2} j_{l}^{2}-\left(j_{i}^{2}-j_{k}^{2}+j_{l}^{2}\right)^{2}\right)}}, \\
\cos \varphi_{i}=\frac{1}{2} \frac{j_{i}^{2}-j_{k}^{2}-j_{l}^{2}-2 m_{k} m_{l}}{\sqrt{\left(j_{k}^{2}-m_{k}^{2}\right)\left(j_{l}^{2}-m_{l}^{2}\right)}}. 
\end{gathered} \label{eq:8:9:16}
\end{equation}
Indices $i, k, l$ in \eqref{eq:8:9:16} are those which may be obtained by cyclic permutation of $1,2,3$. The quantities which enter Eqs.~\eqref{eq:8:9:14}-\eqref{eq:8:9:16} have a simple geometrical interpretation. Three momenta $j_{1}=a+\frac{1}{2}, j_{2}=b+\frac{1}{2}$ and $j_{3}=c+\frac{1}{2}$ form a triangle oriented in space in such a manner that the projections of its sides on the $z$ axis are\\
equal to $m_{1}, m_{2}$ and $m_{3}$, respectively (Fig.~\ref{chap8:fig:1}). $S$ is an area of the triangle projection onto the ( $x, y$ )-plane. Let us construct a trihedral prism whose sidelong edges are parallel to the $z$ axis and one of the cross sections coincides with the oriented triangle. Then $\theta_{i}$ is the angle between the normals to the planes adjacent to the edge $j_{i}$. The angles $\varphi_{i}$ are shown in Fig.~\ref{chap8:fig:1}.

In particular, if $\alpha=\beta=\gamma=0$, we get $\theta_{1}=\theta_{2}=\theta_{3}=\pi / 2$.
\begin{equation}
\clebsch{a}{0}{b}{0}{c}{0} \approx(-1)^{(a+b-c)/2}\left[\frac{1+(-1)^{a+b-c}}{2}\right] \sqrt{\frac{2 c+1}{2 \pi S}} . \label{eq:8:9:17}
\end{equation}
Equations \eqref{eq:8:9:12}-\eqref{eq:8:9:17} are valid, if $S^{2}>0$.

\begin{figure}[tbh]
\begin{center}
  \includegraphics{fig8_1}
\caption{Geometric interpretation of the parameters in the asymptotic relation \eqref{eq:8:9:12}.}
\label{chap8:fig:1}
\end{center}
\end{figure}

More accurate expressions valid not only for $S^{2}>0$ but also if $S^{2}<0$ or $S^{2} \approx 0$ were obtained by Schulten and Gordon \cite{ref140}. In the case of $S^{2} \geq 0$ they derived
\begin{subequations}
\begin{gather}
\clebsch{a}{\alpha}{b}{\beta}{c}{\gamma}=(-1)^{2 b+c+\gamma+1} \sqrt{\frac{2 c+1}{2 S}} Z^{\frac{1}{4}} \notag\\
\times \begin{cases}\cos \Omega_{0} A i(-Z)-\sin \Omega_{0} B i(-Z), & \text { if } \Omega-\Omega_{0}<0, \\
\cos \Omega_{0} B i(-Z)-\sin \Omega_{0} A i(-Z), & \text { if } \Omega-\Omega_{0}>0,\end{cases}, \label{eq:8:9:18a}
\end{gather}
while for $S^{2} \leq 0$
\begin{gather}
\clebsch{a}{\alpha}{b}{\beta}{c}{\gamma}=(-1)^{2 b+c+\gamma+1} \sqrt{\frac{2 c+1}{2|S|}} Z^{\frac{1}{4}} \notag\\
\times \begin{cases}\cos \Omega_{0} A i(Z)-\sin \Omega_{0} B i(Z), & \text { if } \Omega-\Omega_{0}<0, \\
\cos \Omega_{0} B i(Z)-\sin \Omega_{0} A i(Z), & \text { if } \Omega-\Omega_{0}>0 .\end{cases}. \label{eq:8:9:18b}
\end{gather}
\end{subequations}
In these equations $A i(Z)$ and $B i(Z)$ are the Airy functions
\[
\begin{aligned}
Z & =\left(\frac{3}{2}\left|\Omega-\Omega_{0}\right|\right)^{\frac{1}{3}}, \\
\Omega_{0} & =j_{1} \theta_{1}^{(0)}+j_{2} \theta_{2}^{(0)}+j_{3} \theta_{3}^{(0)}+m_{1} \varphi_{2}^{(0)}-m_{2} \varphi_{1}^{(0)}, \\
\Omega & =j_{1} \theta_{1}+j_{2} \theta_{2}+j_{3} \theta_{3}-m_{2} \varphi_{1}+m_{1} \varphi_{2},
\end{aligned}
\]
$\theta_{i}^{(0)}$ and $\varphi_{i}^{(0)}$ being determined by the relations
\[
\begin{aligned}
& \theta_{i}^{(0)}= \begin{cases}0 & \text { if } 0 \leq \operatorname{Re} \theta_{i} \leq \pi / 2 \\
\pi & \text { if } \pi / 2<\operatorname{Re} \theta_{i} \leq \pi\end{cases} \\
& \varphi_{i}^{(0)}= \begin{cases}0 & \text { if } 0 \leq \operatorname{Re} \varphi_{i} \leq \pi / 2 \\
\pi & \text { if } \pi / 2<\operatorname{Re} \varphi_{i} \leq \pi\end{cases}
\end{aligned}
\]
\subsection{Squares of the Clebsch-Gordan Coefficients in the Classical Limit}\label{sec:8.9.4}\index{Clebsch-Gordan coefficient!squares in the classical limit}\index{probability amplitude}
According to Sec.~\ref{sec:8.1.1} the square of a Clebsch-Gordan coefficient represents the probability of the coupling of two angular momenta $\mathrm{j}_{1}$ and $\mathrm{j}_{2}$ with projections $m_{1}$ and $m_{2}$ into the resultant angular momentum j with , projection $m$. On the other hand, the same square is equal to the probability that the momenta $\vect{j}_{1}$ and $\vect{j}_{2}$ coupled into the momentum $j$ with the projection $m$ have the projections $m_{1}$ and $m_{2}$, respectively. Let us consider the square of the Clebsch-Gordan coefficients in the classical limit \cite{ref060}.

If the vector $j$ as well as the magnitudes of $j_{1}$ and $j_{2}$ are given, the possible positions of the end of the vector $\mathrm{j}_{1}$ form a circle (Fig.~\ref{chap8:fig:2}). The probability in question is proportional to the length of the circle arc confined between the planes $z=m_{1}$ and $z=m_{2}$. In the classical limit (for $\left.j_{1}, j_{2}, j \gg 1\right)$ this probability may be expressed by
\begin{equation}
\left(\clebsch{j_{1}}{m_{1}}{j_{2}}{m_{2}}{j}{m}\right)^{2} \approx \frac{2 j+1}{\pi}\left[-\left(j_{1}^{4}+j_{2}^{4}+j^{4}\right)+2\left(j_{1}^{2} j_{2}^{2}+j_{1}^{2} j^{2}+j_{2}^{2} j^{2}\right)-4\left(j_{1}^{2} m_{2} m+j_{2}^{2} m_{1} m-j^{2} m_{1} m_{2}\right)\right]^{-\frac{1}{2}}. \label{eq:8:9:19}
\end{equation}
\eqref{eq:8:9:19} is valid only on the average, because if $j_{1}, j_{2}$ and $j$ are large, even fairly minor relative variations of these quantities lead to rapid oscillations of the Clebsch-Gordan coefficients (See \eqref{eq:8:9:12}). When $m_{1}=m_{2}=m=0$ and $j_{1}+j_{2}-j$ is even, one has
\begin{equation}
\clebsch{j_{1}}{0}{j_{2}}{0}{j}{0} \approx(-1)^{(j_1+j_{2}-j)/2} \sqrt{\frac{2 j+1}{\pi}}\left[-\left(j_{1}^{4}+j_{2}^{4}+j^{4}\right)+2\left(j_{1}^{2} j_{2}^{2}+j_{1}^{2} j^{2}+j_{2}^{2} j^{2}\right)\right]^{-\frac{1}{4}}. \label{eq:8:9:20}
\end{equation}
\begin{figure}[tbh!]
\begin{center}
\includegraphics{fig8_2}
\caption{Addition of angular momenta in the classical limit.}
\label{chap8:fig:2}
\end{center}
\end{figure}

\section{ZEROS OF THE VECTOR-ADDITION COEFFICIENTS}\label{sec:8.10}\index{Clebsch-Gordan coefficient!zeros of}\index{zeros!of the vector-addition coefficients}\index{selection rules}
A Clebsch-Gordan coefficient $\clebsch{a}{\alpha}{b}{\beta}{c}{\gamma}$ and associated $3 j m$ symbol $\threej{a}{b}{c}{\alpha}{\beta}{-\gamma}$ vanish unless triangular conditions are satisfied (Eqs.~\eqref{eq:8:1:1} and \eqref{eq:8:1:2}). However, there exist non-trivial cases when the Clebsch-Gordan coefficients are equal to zero although the triangular conditions are fulfilled. This leads to some additional selection rules which forbid the quantum transitions whose amplitudes are proportional to the appropriate Clebsch-Gordan coefficients. Here we point out the most important relations between the arguments for which the Clebsch-Gordan coefficients vanish.
\begin{enumerate}[label=(\alph*)]
\item $\clebsch{a}{0}{b}{0}{c}{0}=0$, if $J \equiv a+b+c$ is odd;
\item $\clebsch{a}{\alpha}{a}{\alpha}{c}{\gamma}=0$, $\clebsch{a}{-\alpha}{c}{\gamma}{a}{\alpha}=0$, $\clebsch{c}{\gamma}{a}{-\alpha}{a}{\alpha}=0$, if $J=2 a+c$ is odd;
\item $\clebsch{a}{J-3 a}{b}{J-3 b}{c}{-(J-3 c)}=0$, if $J$ is odd;
\item $\clebsch{a}{\alpha}{b}{\beta}{a+b-1}{\gamma}=0$, if $\frac{\alpha}{\beta}=\frac{a}{b}$;
\item $\clebsch{a}{\alpha}{b}{\beta}{a-b+1}{\gamma}=0$ if $\frac{\alpha}{\beta}=-\frac{a+1}{b}$;
\item $\clebsch{a}{\alpha}{b}{\beta}{c}{c-1}=0$ if $a(a+1)-b(b+1)=(\alpha-\beta) c$, in particular
\[
\clebsch{a}{\alpha}{a}{\alpha}{c}{c-1}=0 .
\]
\item $\clebsch{a}{\alpha}{b}{b-1}{c}{\gamma}=0$ if $a(a+1)-c(c+1)=-(\alpha+\gamma) b$ in particular,
\[
\clebsch{a}{\alpha}{b}{b-1}{a}{-\alpha}=0 .
\]
\end{enumerate}
Below we list particular 3jm symbols with $J=a+b+c \leq 17$ which vanish in nontrivial cases. They are arranged in the order of increasing $J$ and $a \leq b \leq c$. Roots of other 3jm symbols may be obtained by the use of the simplest symmetry properties of the $3 j \mathrm{~m}$ symbols, concerned with the permutations of momenta $a, b, c$ and sign changes of the projections $\alpha, \beta, \gamma$ (see Eqs.~\eqref{eq:8:4:5} and \eqref{eq:8:4:6}), Hence, any $3 j m$ symbols from the table given below generally represents 12 formally different versions of the $3 j m$ symbols or the Clebsch-Gordan coefficients which are equal to zero. One should bear in mind that the $3 j \mathrm{~m}$ symbols and the Clebsch-Gordan coefficients are related by (see Sec.~\ref{sec:8.1.2})
\[
\threej{a}{b}{c}{\alpha}{\beta}{-\gamma}=(-1)^{a+b+\gamma} \frac{1}{\sqrt{2 c+1}} \clebsch{a}{\alpha}{b}{\beta}{c}{\gamma} .
\]
\section{CONNECTION OF THE CLEBSCH-GORDAN COEFFICIENTS AND THE 3 jm SYMBOLS WITH ANALOGOUS FUNCTIONS OF OTHER AUTHORS}\label{sec:8.11}\index{Clebsch-Gordan coefficient!notations of other authors}\index{notation!of other authors}
\[
\clebsch{j_{1}}{m_{1}}{j_{2}}{m_{2}}{j}{m}= \begin{cases}S_{j m_{1} m_{2}}^{j j_{2}} & \text { - Wigner \cite{ref043}, } \\
A_{m m_{1} m_{2}}^{j j_{1} j_{2}} & \text { - Eckart \cite{ref015}, } \\
C_{m_{1} m_{2}}^{j} & \text { - Vander Waerden \cite{ref040}, } \\
\left(j_{1} j_{2} m_{1} m_{2} \mid j_{1} j_{2} j m\right) & \text { - Condon and Shortley \cite{ref010}, } \\
C_{j m}^{j_{1} j_{2}}\left(m_{1} m_{2}\right) & \text { - Fock \cite{ref068}, } \\
X\left(j, m_{1}, j_{1}, j_{2}, m_{1}\right) & \text { - Boys \cite{ref059}, } \\
C\left(j m_{1} m_{1} m_{2}\right) & \text { - Blatt and Weisskopf \cite{ref007}, } \\
C_{m_{1} m_{2} m}^{j_{1} j_{2} j} & \text { - Biedenharn \cite{ref115}, } \\
C\left(j_{1} j_{2} j, m_{1} m_{2}\right) & \text { - Rose \cite{ref030}, } \\
{\left[\begin{array}{lll}
j_{1} & j_{2} & j \\
m_{1} & m_{2} & m
\end{array}\right]} & \text { - Yutsis and Bandzaitis \cite{ref045}, } \\
\braket{j_{1} m_{1} j_{2} m_{2}}{\left(j_{1} j_{2}\right) j m} & \text { - Fano \cite{ref066}. }\end{cases}
\]
\begin{table}[tbhp!]
  \caption{Accidental zeroes of the $3j$ coefficients for $J=j_1+j_2+j_3\leq 17$}
$\small\setlength{\arraycolsep}{2pt}
\begin{aligned}
 J=3& \threej{1}{1}{1}{0}{0}{0} . \\
J=5 &\threej{1}{2}{2}{0}{0}{0},\threej{3 / 2}{3 / 2}{2}{1 / 2}{1 / 2}{-1} . \\
J=7 &\threej{1}{3}{3}{0}{0}{0},\threej{2}{2}{3}{0}{0}{0},\threej{2}{2}{3}{1}{1}{-2},\threej{2}{5 / 2}{5 / 2}{-1}{1 / 2}{1 / 2} . \\
J=8 &\threej{2}{3}{3}{0}{2}{-2},\threej{3 / 2}{3}{7 / 2}{1 / 2}{1}{-3 / 2} \text {. } \\
J=9 &\threej{1}{4}{4}{0}{0}{0},\threej{2}{3}{4}{0}{0}{0},\threej{3}{3}{3}{0}{0}{0},\threej{3}{3}{3}{1}{1}{-2}, \\
& \threej{2}{7 / 2}{7 / 2}{-1}{1 / 2}{1 / 2},\threej{5 / 2}{3}{7 / 2}{3 / 2}{0}{-3 / 2},\threej{5 / 2}{5 / 2}{4}{1 / 2}{1 / 2}{-1},\threej{5 / 2}{5 / 2}{4}{3 / 2}{3 / 2}{-3} . \\
J=11 &\threej{1}{5}{5}{0}{0}{0},\threej{2}{4}{5}{0}{0}{0},\threej{2}{4}{5}{1}{2}{-3},\threej{3}{3}{5}{0}{0}{0},\threej{3}{3}{5}{1}{1}{-2},\threej{3}{3}{5}{2}{2}{-4}, \\
& \threej{3}{4}{4}{0}{0}{0},\threej{3}{4}{4}{-2}{1}{1},\threej{2}{9 / 2}{9 / 2}{-1}{1 / 2}{1 / 2},\threej{3}{7 / 2}{9 / 2}{2}{1 / 2}{-5 / 2},\threej{7 / 2}{7 / 2}{4}{1 / 2}{1 / 2}{-1}, \\
& \threej{7 / 2}{7 / 2}{4}{3 / 2}{3 / 2}{-3},\threej{3 / 2}{9 / 2}{5}{1 / 2}{3 / 2}{-2},\threej{5 / 2}{4}{9 / 2}{1 / 2}{3}{-7 / 2} \text {. } \\
J=13 &\threej{1}{6}{6}{0}{0}{0},\threej{2}{5}{6}{0}{0}{0},\threej{3}{4}{6}{0}{0}{0},\threej{3}{5}{5}{0}{0}{0},\threej{3}{5}{5}{-2}{1}{1},\threej{4}{4}{5}{0}{0}{0}, \\
& \threej{4}{4}{5}{1}{1}{-2},\threej{4}{4}{5}{2}{2}{-4},\threej{2}{11 / 2}{11 / 2}{-1}{1 / 2}{1 / 2},\threej{7 / 2}{9 / 2}{5}{5 / 2}{-1 / 2}{-2},\threej{7 / 2}{4}{11 / 2}{5 / 2}{1}{-7 / 2}, \\
& \threej{4}{9 / 2}{9 / 2}{-1}{1 / 2}{1 / 2},\threej{4}{9 / 2}{9 / 2}{-3}{3 / 2}{3 / 2},\threej{7 / 2}{7 / 2}{6}{1 / 2}{1 / 2}{-1},\threej{7 / 2}{7 / 2}{6}{3 / 2}{3 / 2}{-3},\threej{7 / 2}{7 / 2}{6}{5 / 2}{5 / 2}{-5} . \\
J=14 &\threej{3}{5}{6}{1}{-1}{0},\threej{3}{5}{6}{2}{1}{-3},\threej{3}{5}{6}{1}{4}{-5},\threej{4}{4}{6}{0}{2}{-2},\threej{4}{5}{5}{1}{3}{-4},\threej{4}{5}{5}{2}{1}{-3}, \\
& \threej{3}{9 / 2}{13 / 2}{1}{3 / 2}{-5 / 2},\threej{5 / 2}{5}{13 / 2}{1 / 2}{1}{-3 / 2},\threej{5 / 2}{5}{13 / 2}{3 / 2}{3}{-9 / 2},\threej{4}{9 / 2}{11 / 2}{0}{7 / 2}{-7 / 2},\threej{3 / 2}{6}{13 / 2}{1 / 2}{2}{-5 / 2},\threej{5 / 2}{11 / 2}{6}{3 / 2}{1 / 2}{-2} . \\
J=15 &\threej{1}{7}{7}{0}{0}{0},\threej{2}{6}{7}{0}{0}{0},\threej{2}{6}{7}{1}{3}{-4},\threej{3}{5}{7}{0}{0}{0},\threej{3}{6}{6}{0}{0}{0},\threej{3}{6}{6}{0}{5}{-5},\threej{3}{6}{6}{-2}{1}{1}, \\
& \threej{4}{5}{6}{0}{0}{0},\threej{4}{5}{6}{3}{0}{-3},\threej{4}{4}{7}{0}{0}{0},\threej{4}{4}{7}{1}{1}{-2},\threej{4}{4}{7}{2}{2}{-4},\threej{4}{4}{7}{3}{3}{-6},\threej{5}{5}{5}{0}{0}{0}, \\
& \threej{5}{5}{5}{1}{1}{-2},\threej{5}{5}{5}{2}{2}{-4},\threej{2}{13 / 2}{13 / 2}{-1}{1 / 2}{1 / 2},\threej{4}{9 / 2}{13 / 2}{3}{3 / 2}{-9 / 2},\threej{4}{11 / 2}{11 / 2}{-1}{1 / 2}{1 / 2},\threej{4}{11 / 2}{11 / 2}{-3}{3 / 2}{3 / 2}, \\
& \threej{9 / 2}{5}{11 / 2}{3 / 2}{0}{-3 / 2},\threej{9 / 2}{9 / 2}{6}{1 / 2}{1 / 2}{-1},\threej{9 / 2}{9 / 2}{6}{3 / 2}{3 / 2}{-3},\threej{9 / 2}{9 / 2}{6}{5 / 2}{5 / 2}{-5} . \\
J=17 &\threej{1}{8}{8}{0}{0}{0},\threej{2}{7}{8}{0}{0}{0},\threej{3}{6}{8}{0}{0}{0},\threej{3}{6}{8}{2}{4}{-6},\threej{3}{7}{7}{0}{0}{0},\threej{3}{7}{7}{-2}{1}{1},\threej{4}{5}{8}{0}{0}{0},\threej{4}{6}{7}{0}{0}{0}, \\
& \threej{5}{5}{7}{0}{0}{0},\threej{5}{5}{7}{1}{1}{-2},\threej{5}{5}{7}{2}{2}{-4},\threej{5}{5}{7}{3}{3}{-6},\threej{5}{6}{6}{0}{0}{0},\threej{5}{6}{6}{-2}{1}{1},\threej{5}{6}{6}{-4}{2}{2}, \\
& \threej{2}{15 / 2}{15 / 2}{-1}{1 / 2}{1 / 2},\threej{3}{13 / 2}{15 / 2}{2}{3 / 2}{-7 / 2},\threej{4}{13 / 2}{13 / 2}{-1}{1 / 2}{1 / 2},\threej{4}{13 / 2}{13 / 2}{-3}{3 / 2}{3 / 2},\threej{5}{11 / 2}{13 / 2}{2}{1 / 2}{-5 / 2},\threej{7 / 2}{6}{15 / 2}{3 / 2}{5}{-13 / 2}, \\
& \threej{9 / 2}{6}{13 / 2}{1 / 2}{-5}{9 / 2},\threej{9 / 2}{6}{13 / 2}{-7 / 2}{1}{5 / 2},\threej{11 / 2}{11 / 2}{6}{1 / 2}{1 / 2}{-1},\threej{11 / 2}{11 / 2}{6}{3 / 2}{3 / 2}{-3},\threej{11 / 2}{11 / 2}{6}{5 / 2}{5 / 2}{-5},\threej{9 / 2}{11 / 2}{7}{7 / 2}{1 / 2}{-4}, \\
& \threej{3 / 2}{15 / 2}{8}{1 / 2}{5 / 2}{-3},\threej{9 / 2}{9 / 2}{8}{1 / 2}{1 / 2}{-1},\threej{9 / 2}{9 / 2}{8}{3 / 2}{3 / 2}{-3},\threej{9 / 2}{9 / 2}{8}{5 / 2}{5 / 2}{-5},\threej{9 / 2}{9 / 2}{8}{7 / 2}{7 / 2}{-7},\threej{9 / 2}{5}{15 / 2}{7 / 2}{2}{-11 / 2} .
\end{aligned}
$
\end{table}
\[
\threej{j_{1}}{j_{2}}{j_{3}}{m_{1}}{m_{2}}{m_{3}}= \begin{cases}(-1)^{-j_{1}+j_{2}+j_{3}} V\left(j_{1} j_{2} j_{3}, m_{1} m_{2} m_{3}\right) & - \text { Racah \cite{ref091}} \\
(-1)^{-j_{1}+j_{2}+j_{3}} \bar{V}\left(j_{1} j_{2} j_{3}, m_{1} m_{2} m_{3}\right) & - \text { Fano and Racah \cite{ref018} } \\
(-1)^{j_{1}-j_{2}+j_{3}} S_{j_{1} m_{1} j_{2} m_{2} j_{3} m_{3}} & - \text { Landau and Lifshitz \cite{ref025} } \\
(-1)^{j_{1}-j_{2}+j_{3}}\left(j_{1} m_{1}, j_{2} m_{2}, j_{3} m_{3}|0\rangle\right. & - \text { Fano \cite{ref066} } \\
X\left(j_{1} j_{2} j_{3}, m_{1} m_{2} m_{3}\right) & - \text { Schwinger \cite{ref101} } \\
(-1)^{j_{1}-j_{2}+j_{3}} U\left(j_{1} j_{2} m_{1} m_{2} \mid j_{3}-m_{3}\right) & - \text { Lubarskii \cite{ref026} } \\
\left(j_{1} j_{2} j_{3}\right)_{m_{1} m_{2} m_{3}} & - \text { Sharp \cite{ref102} }\end{cases}
\]
\section{ALGEBRAIC TABLES OF THE CLEBSCH-GORDAN COEFFICIENTS}\label{sec:8.12}\index{Clebsch-Gordan coefficient!algebraic tables}\index{tables!of the Clebsch-Gordan coefficients}
Below we present tables of algebraic formulas for the Clebsch-Gordan coefficients $\clebsch{a}{\alpha}{b}{\beta}{c}{\gamma}$ with $b=1 / 2,1,3 / 2,2$, $5 / 2,3,7 / 2,4,9 / 2,5$. If $b=3,7 / 2,4,9 / 2,5$, we give $\clebsch{a}{\alpha}{b}{\beta}{c}{\gamma}$ with $\beta \geq 0$, while $\clebsch{a}{\alpha}{b}{\beta}{c}{\gamma}$ with $\beta<0$ may be obtained from the symmetry property
\[
\clebsch{a}{\alpha}{b}{\beta}{c}{\gamma}=(-1)^{a+b-c} \clebsch{a}{-\alpha}{b}{-\beta}{c}{-\gamma} .
\]
The algebraic tables of the Clebsch-Gordan coefficients are also available in Ref.~\cite{ref010} for $b=1 / 2,1,3 / 2,2$; in Ref.~$[121,125]$ for $b=5 / 2$; in Refs.~\cite{ref117,ref130} with $b=3$; and in Ref.~\cite{ref126} with $b=7 / 2,4,9 / 2,5$.

\section{NUMERICAL TABLES OF THE CLEBSCH-GORDAN COEFFICIENTS}\label{sec:8.13}\index{Clebsch-Gordan coefficient!numerical tables}\index{tables!numerical}
In Table 8.11 we present numerical values of the Clebsch-Gordan coefficients for $a, b, c \leq 3$. These values are given in the form of square roots of rational fractions and in decimals. The Clebsch-Gordan coefficients are separated into groups with respect to the values of $c$. The arguments $a, \alpha, b, \beta$ for the Clebsch-Gordan coefficients given in Table 8.11 satisfy the inequalities (a) $a \geq b, \alpha \geq 0$; (b) $\alpha \geq \beta$ for $a=b$. Other ClebschGordan coefficients with $a, b, c \leq 3$ may be reduced to those in Table 8.11 by the use of the symmetry properties
\[
\clebsch{a}{\alpha}{b}{\beta}{c}{\alpha+\beta}=(-1)^{a+b-c} \clebsch{b}{\beta}{a}{\alpha}{c}{\alpha+\beta}=(-1)^{a+b-c} \clebsch{a}{-\alpha}{b}{-\beta}{c}{-\alpha-\beta}=\clebsch{b}{-\beta}{a}{-\alpha}{c}{-\alpha-\beta}
\]
Numerical values of the Clebsch-Gordan coefficients are also given in Refs.~\cite{ref126, ref127}. Reference \cite{ref127} presents the Clebsch-Gordan coefficients for $\frac{1}{2} \leq a \leq 4, \frac{1}{2} \leq b, c \leq 9 / 2$ in decimals. Reference \cite{ref126} gives the Clebsch-Gordan coefficients for $\frac{1}{2} \leq b \leq 6, a=5,11 / 2,6$ in the form of rational fractions. Numerical values of the $3 j m$ symbols are also available in Ref.~\cite{ref113}.

Tables~\ref{chap8:tab:1}--\ref{chap8:tab:2}. Algebraic Formulas for the Clebsch-Gordan Coefficients.
Tables not reproduced here can be found in the original book.
\begin{table}[tbh!]
\begin{center}
\caption{$\clebsch{a}{\alpha}{\frac{1}{2}}{\beta}{c}{\gamma}$}
\label{chap8:tab:1}
\begin{tabular}{c|c|c}
\hline
$c$ & $\beta=1 / 2$ & $\beta=-1 / 2$ \\
\hline
$a+1 / 2$ & $\left[\frac{c+\gamma}{2 c}\right]^{1 / 2}$ & $\left[\frac{c-\gamma}{2 c}\right]^{1 / 2}$ \\
$a-1 / 2$ & $-\left[\frac{c-\gamma+1}{2 c+2}\right]^{1 / 2}$ & $\left[\frac{c+\gamma+1}{2 c+2}\right]^{1 / 2}$ \\
\end{tabular}
\end{center}
\end{table}

\begin{table}[tbh!]
\begin{center}
\caption{$\clebsch{a}{\alpha}{1}{\beta}{c}{\gamma}$}
\label{chap8:tab:2}
\begin{tabular}{c|c|c|c}
\hline
c & $\beta=1$ & $\beta=0$ & $\beta=-1$ \\
\hline
$a+1$ & $\left[\frac{(c+\gamma-1)(c+\gamma)}{(2 c-1) 2 c}\right]^{1 / 2}$ & 
\(\left[\frac{(c+\gamma)(c-\gamma)}{(2 c-1) c}\right]^{1 / 2}\)& 
\(\left[\frac{(c-\gamma-1)(c-\gamma)}{(2 c-1) 2c}\right]^{1 / 2}\)
\\
$a$ & \(-\left[\frac{(c+\gamma)(c-\gamma+1)}{2 c(c+1)}\right]^{1 / 2 }\) &\(\frac{\gamma}{[c(c+1)]^{1 / 2}}\)&
$\left[\frac{(c+\gamma+1)(c-\gamma)}{2 c(c+1)}\right]^{1 / 2}$\\
$a-1$ & \(  \left[\frac{(c-\gamma+1)(c-\gamma+2)}{(2 c+2)(2 c+3)}\right]^{1 / 2} \) 
 & \(-\left[\frac{(c+\gamma+1)(c-\gamma+1)}{( c+1)(2 c+3)}\right]^{1 / 2} \)
 & \( \left[\frac{(c+\gamma+2)(c+\gamma+1)}{(2 c+2)(2 c+3)}\right]^{1 / 2} \) 
\end{tabular}
\end{center}
\end{table}

